%% Supplementary Materials of the Crystals (MDPI) article in paper/main.tex. One PDF, sections S1 to S9,
%% figures S1 to S8 (final numbering: docs/figure_numbering.md), tables S1 onwards.
%% Build from paper/si with `latexmk` (paper/si/latexmkrc) after `make -C paper numbers`; every number in the
%% prose is a macro of paper/generated/numbers.tex (rules: paper/tools/README.md, docs/prose_style.md).
%% Tables: sections/tab_*.tex and tables/*.tex are written by tools/make_si_tables.py.
\documentclass[10pt]{article}
\usepackage{sisetup}

\input{../generated/numbers}

\hypersetup{pdftitle={Supplementary Materials: Oxygen-vacancy formation energies in non-magnetic oxides},
  pdfauthor={Palaash Gang and Aiden Lee}}

\begin{document}

\begin{center}
{\large\bfseries Supplementary Materials}\\[6pt]
{\bfseries Oxygen-vacancy formation energies in non-magnetic oxides: how much do explicit electronic descriptors
add to graph networks as training data grows?}\\[6pt]
Palaash Gang and Aiden Lee
\end{center}

\vspace{6pt}
\noindent\begin{minipage}{\linewidth}\small
These Supplementary Materials accompany the article in \textit{Crystals}. Sections~\ref{si:data} to
\ref{si:repro} give the data filter and descriptor classes, the training protocol and tuned settings, the split
designs, the additional within-host analyses, the hardware variation, the sensitivity analyses, the latent probe,
the structure renders and the reproducibility record. Figures~S1 to S8 and the tables are numbered with the
prefix S; figure and table numbers in the article refer to the article. Every result carries the label of
its provenance: registered (the originally registered network, v1, with a rule fixed before its results), post
hoc (the primary network, v2, adopted after diagnostics of v1, or a rule fixed after its results) or extension
(analysis beyond the registered plan).
\end{minipage}

\section{Data: filter cascade, descriptor classes and integrity exclusions}
\label{si:data}

The dataset of this work is the set of neutral oxygen vacancies, $V_\mathrm{O}^{\times}$, of the release of
Kumagai et al.\ \citep{Kumagai2021insights} that survive a fixed cascade of filters (Table~\ref{tab:cascade}; the
cascade is drawn as Figure~\ref{fig:cascade} in Section~\ref{si:splits}). The release holds
\numFilterCascadeStepsReleaseSites{} completed neutral calculations in \numFilterCascadeStepsReleaseHosts{} host
oxides. Kumagai et al.\ built from them a neutral machine-learning subset of \numFilterCascadeStepsMlSubsetSites{}
vacancies in \numFilterCascadeStepsMlSubsetHosts{} hosts, from which they had already removed the
perturbed-host-state entries (\numScreeningAbsenceReleaseNeutralCompletedPhs{} of the completed neutral entries),
the entries whose band edges could not be determined and the dynamically unstable hosts. A perturbed-host-state
entry is a vacancy whose two donor electrons do not stay in a localised in-gap state but delocalise into the host
conduction band, so the subset, and every result of this work, describes deep, localised neutral vacancies in
ordered oxides classified as non-magnetic in the release. Our three further filters, D1 to D3, remove what the data audit
(\texttt{docs/data\_audit.md}) found inconsistent in that subset; they leave \numFilterCascadeStepsFinalSites{}
oxygen sites in \numFilterCascadeStepsFinalHosts{} hosts, one host per formula. The site counts after each
step are read from \texttt{results/filter\_cascade.json}.

\begin{table}[htbp]
\caption{Filter cascade of the \numFilterCascadeStepsFinalHosts{}-host dataset, with the oxygen sites, hosts and formulas left after each step
(\texttt{results/filter\_cascade.json}; the same table as \arttabcascade{} of the article).}
\label{tab:cascade}
\centering\small
\begin{adjustbox}{max width=\linewidth}
\begin{tabular}{llrrr}
\toprule
Step & Criterion & Sites & Hosts & Formulas \\
\midrule
Release & Completed neutral entries in the release & 2103 & 935 & 935 \\
ML subset & Kumagai neutral ML subset (\texttt{charge0.csv}) & 1745 & 824 & 824 \\
D1 & Vacancy split, unknown type or not same config from init & 1730 & 819 & 819 \\
D2 & Host Sr$_{2}$Zr$_{7}$O$_{16}$ (ML site labels do not match release targets) & 1726 & 818 & 818 \\
D3 & Ba$_{4}$Nb$_{2}$WO$_{12}$, Sr$_{2}$SnO$_{4}$: vacancy site is O with matching first cation shell & 1726 & 818 & 818 \\
\midrule
Final & Final universe & 1726 & 818 & 818 \\
\bottomrule
\end{tabular}

\end{adjustbox}
\end{table}

The three exclusions concern the integrity of the labels, not the chemistry of the hosts. D1 removes the
\numFilterCascadeStepsDOneNRemoved{} entries of the subset that the release itself flags as a split vacancy, an
unknown defect type, or a relaxed configuration that is not the configuration of the initial vacancy. Each of
them also carries the last flag, so the stored energy of these entries may not belong to a single vacancy on the
labelled site, which is the structure that a model reads. The subset had kept these entries although the release
flags them. D2 removes the
\numFilterCascadeStepsDTwoNRemoved{} sites of Sr$_2$Zr$_7$O$_{16}$, whose site labels in the machine-learning
tables do not match the targets of the release; the release gives no criterion that explains the mismatch, so the
host is excluded rather than repaired. D3 concerns the two hosts whose cell metadata are inconsistent with their
supercells, Ba$_4$Nb$_2$WO$_{12}$ and Sr$_2$SnO$_4$ (the Sr$_2$SnO$_4$ supercell of
\numStructureAuditAnomalousHostsSrTwoSnOFourNAtomsSc{} atoms has \numStructureAuditAnomalousHostsSrTwoSnOFourVolumeRatio{}
of the expected volume, because it was built from the primitive cell). An entry of these hosts is kept only if
its labelled vacancy site maps to an oxygen atom of the supercell whose first cation shell matches that of the
unit-cell site within \numFilterCascadeDThreeSrTwoSnOFourSitesOOneToleranceA{}~\AA{}; every entry passed, and D3
removed \numFilterCascadeStepsDThreeNRemoved{} sites. Every site of the dataset lies in the published split of
Kiyohara et al.\ \citep{Kiyohara2025machinelearning} (\numFilterCascadeKiyoharaSplitSitesAbsent{} sites absent),
which is used for the comparison with the published accuracy (Section~\ref{si:splits}).

The descriptors of the release fall into three classes (Table~\ref{tab:taxonomy}). The host-electronic class
holds properties of the pristine host from its DFT electronic structure: the band gap of the PBEsol(+$U$)
calculation, the electronic and ionic dielectric constants and the orbital character of the band edges, together
with the formation energy of the host oxide, which measures how strongly oxygen is bound. The site-electronic class
describes the oxygen atom to be removed in the pristine host: its Bader charge and volume, its Born effective
charge, the same quantities averaged over and extremal among its Voronoi neighbours, and the centre of its
O-2$p$ band below the valence-band maximum. The structural and compositional class holds the fractions of the
solid angle of the oxygen Voronoi cell that face each element and the electronegativities of the neighbours. D
receives the two electronic classes, \numRfElectronicNFeatures{} descriptors; RF-Kumagai uses all
\numRfKumagaiNFeatures{} descriptors, RF-electronic the \numRfElectronicNFeatures{} electronic ones and
RF-structural the \numRfStructuralNFeatures{} structural and compositional ones. Every descriptor is computed on
the pristine host, so none of them contains the relaxed defect geometry or the defect energy. The host formation
energy is named \texttt{formation\_energy} in the release and is not the vacancy formation energy that the models
predict.

{\footnotesize\tighttabs
% Generated by paper/si/tools/make_si_tables.py from docs/descriptor_taxonomy.csv. Do not edit.
\begin{longtable}{@{}>{\raggedright}p{0.25\linewidth}>{\raggedright}p{0.15\linewidth}>{\raggedright}p{0.135\linewidth}>{\raggedright\arraybackslash}p{0.385\linewidth}@{}}
\caption{The descriptors of the release of Kumagai et al.\ \citep{Kumagai2021insights} and their classes, transcribed from \texttt{docs/descriptor\_taxonomy.csv} (names as in the release; definition abridged to its first sentence). The host-electronic and site-electronic DFT descriptors are those given to D; RF-Kumagai uses all of them, RF-electronic the two electronic classes and RF-structural the structural and compositional class. Every descriptor is a property of the pristine host; \texttt{formation\_energy} is the formation energy of the host oxide, not the vacancy target.}\label{tab:taxonomy}\\
\toprule
Name & Class & Units & Definition \\
\midrule
\endfirsthead
\toprule
Name & Class & Units & Definition \\
\midrule
\endhead
\bottomrule
\endlastfoot
\texttt{ave\_ele\_diele} & host-electronic & dimensionless & Spherically averaged electronic (ion-clamped) dielectric constant of the pristine host, from DFPT \\
\texttt{ave\_ion\_diele} & host-electronic & dimensionless & Spherically averaged ionic dielectric constant of the pristine host, from DFPT \\
\texttt{band\_gap} & host-electronic & eV & Band gap of the pristine host from the PBEsol(+U) calculation; not the nsc-dd hybrid gap \\
\texttt{cbm\_d} & host-electronic & dimensionless & Fraction of d-orbital character in the pristine host's conduction-band minimum state (projection ratio) \\
\texttt{cbm\_f} & host-electronic & dimensionless & Fraction of f-orbital character in the pristine host's conduction-band minimum state (projection ratio) \\
\texttt{cbm\_p} & host-electronic & dimensionless & Fraction of p-orbital character in the pristine host's conduction-band minimum state (projection ratio) \\
\texttt{cbm\_s} & host-electronic & dimensionless & Fraction of s-orbital character in the pristine host's conduction-band minimum state (projection ratio) \\
\texttt{formation\_energy} & host-electronic & eV/atom & Host formation energy per atom retrieved from the Materials Project database (empirically corrected) \\
\texttt{vbm\_d} & host-electronic & dimensionless & Fraction of d-orbital character in the pristine host's valence-band maximum state (projection ratio) \\
\texttt{vbm\_f} & host-electronic & dimensionless & Fraction of f-orbital character in the pristine host's valence-band maximum state (projection ratio) \\
\texttt{vbm\_p} & host-electronic & dimensionless & Fraction of p-orbital character in the pristine host's valence-band maximum state (projection ratio) \\
\texttt{vbm\_s} & host-electronic & dimensionless & Fraction of s-orbital character in the pristine host's valence-band maximum state (projection ratio) \\
\addlinespace
\texttt{ave\_bec} & site-electronic & e & Spherically averaged Born effective charge (trace/3) of the O atom to be removed, in the pristine host (DFPT) \\
\texttt{bader\_charge} & site-electronic & e & Bader charge of the O atom to be removed, from the pristine host unit-cell all-electron density; net ionic charge, negative for the anion \\
\texttt{bader\_volume} & site-electronic & \AA$^{3}$ & Bader (zero-flux) volume of the O atom to be removed, from the pristine host all-electron density \\
\texttt{nn\_ave\_ave\_bec} & site-electronic & e & Solid-angle-weighted mean of the spherically averaged Born effective charge over the Voronoi neighbours \\
\texttt{nn\_ave\_bader\_charge} & site-electronic & e & Solid-angle-weighted mean of the Bader charge over the Voronoi neighbours of the O site \\
\texttt{nn\_max\_ave\_bec} & site-electronic & e & Maximum of the spherically averaged Born effective charge over the Voronoi neighbours \\
\texttt{nn\_max\_bader\_charge} & site-electronic & e & Maximum of the Bader charge over the Voronoi neighbours of the O site \\
\texttt{nn\_min\_ave\_bec} & site-electronic & e & Minimum of the spherically averaged Born effective charge over the Voronoi neighbours \\
\texttt{nn\_min\_bader\_charge} & site-electronic & e & Minimum of the Bader charge over the Voronoi neighbours of the O site \\
\texttt{o2p\_center\_from\_vbm} & site-electronic & eV & Centroid of the O-2p projected density of states of the O site, measured as its distance below the VBM (Deml et al \\
\addlinespace
\texttt{nn\_ave\_eleneg} & structural & Pauling units & Solid-angle-weighted mean of the Pauling electronegativity over the Voronoi neighbours (cations and O) \\
\texttt{nn\_max\_eleneg} & structural & Pauling units & Maximum of the Pauling electronegativity over the Voronoi neighbours (cations and O) \\
\texttt{nn\_min\_eleneg} & structural & Pauling units & Minimum of the Pauling electronegativity over the Voronoi neighbours (cations and O) \\
\texttt{w\_X} & structural & fraction & Fraction of the solid angle of the O-site Voronoi cell (pristine unit cell) that faces element X, for X in: Bi, Ag, O, Hg, Au, Ge, Al, Pb, Zn, As, Cu, Sb, Mo, Tl, B, Ba, Ca, Pd, Cd, Pt, Ce, Hf, In, Ga, La, Ta, Mg, Si, Sc, Sn, Ti, Y, Zr, Nb, Na, Sr, V, Li, Be, W, Cr, Cs, K, Re, Rb \\
\end{longtable}
}

\section{Models and training}
\label{si:training}

\subsection{Readout of the structure-only network}

The structure-only network S maps the pristine supercell, with the oxygen atom to be removed marked, to the
formation energy of the neutral vacancy (Figure~\ref{fig:readout}). Successive MEGNet blocks
\citep{Chen2019graph} aggregate information over growing coordination shells around every atom, starting from the
learned element embeddings, the vacancy flag, a Gaussian expansion of each interatomic distance and a global state
of the cell. The final step joins three learned vectors: the description of the vacancy site, a pooled description
of all atoms of the host (Set2Set or the mean, a tuned choice) and the final global state; a small multilayer
network maps them to one energy. The split follows the variance of the target. The pooled host vector and the
global state carry the host-level thermodynamic term, the cost of releasing an oxygen atom from that oxide, and
the vacancy-site vector carries the local environment that separates the sites of one host. In D (the D-state
injection), the host-electronic descriptors enter the initial global state and the site-electronic descriptors the
initial description of the vacancy site, each standardised with training-set statistics, so that each class enters
where its physics acts; the readout itself is unchanged. The readout is the same for both networks, the
originally registered network (v1) and the primary network (v2).

\begin{figure}[htbp]
\centering
\sifig{fig_readout}
\caption{Readout of the structure-only network S, from the initial features of the vacancy site \xvzero{}, the other atoms
\xizero{}, the bonds \eijzero{} and the cell \uzero{}, through $n$ MEGNet blocks, to the vacancy-site vector
$\mathbf{h}_\mathrm{v}$, the pooled vector $\bar{\mathbf{h}}$ and the global state $u$ that a multilayer perceptron
maps to $E_\mathrm{f}$. Green arrows mark where D (D-state) receives the host-electronic descriptors, at \uzero{},
and the site-electronic descriptors, at \xvzero{}.}
\label{fig:readout}
\end{figure}

\subsection{Training protocol, search space and tuned values}

Every network was trained on the standardised formation energy with an absolute-error loss and the AdamW
optimiser; the learning rate was halved (factor \numCfgTrainingSchedulerFactor{}) when the validation error had
not improved for \numCfgTrainingSchedulerPatience{} passes over the training hosts. The validation hosts that
decide early stopping are carved out of the training hosts of each budget (a fraction of
\numCfgTrainingValidationFrac{}, at least \numCfgTrainingValidationMinHosts{} hosts), so a budget counts every
host whose DFT label a model consumes. Training stops when the validation error has not improved for
\numCfgTrainingEarlyStoppingPatience{} passes, or at the cap of \numCfgTrainingMaxEpochs{} passes, and the weights
of the best pass are restored.

The settings of each network were tuned once, before any test host was evaluated, on the nested training sets of
resample \numCfgTuningAnchorResample{} at three anchor budgets, \numCfgTuningAnchorsItemOne{},
\numCfgTuningAnchorsItemTwo{} and \numCfgTuningAnchorsItemThree{} hosts (\texttt{docs/tuning.md}). Each model and
anchor had one study of \numCfgTuningOptunaTrialsPerModelPerAnchor{} complete trials with a tree-structured
Parzen sampler, scored by the best validation error; at the largest anchor, \numTuningSummaryStudiesSASixHundredFiftyFourValidationNTrainHosts{}
hosts were fitted and \numTuningSummaryStudiesSASixHundredFiftyFourValidationNValHosts{} validated. The search
space was the same for every model (Table~\ref{tab:space}). A training budget uses the values of its anchor: the
budgets \numCfgBudgetsHostsItemOne{} and \numCfgBudgetsHostsItemTwo{} those of anchor
\numCfgTuningBudgetToAnchorTwentyFive{}, the budgets \numCfgBudgetsHostsItemThree{} and
\numCfgBudgetsHostsItemFour{} those of anchor \numCfgTuningBudgetToAnchorOneHundred{}, and the budgets
\numCfgBudgetsHostsItemFive{} and \numCfgBudgetsHostsItemSix{} those of anchor
\numCfgTuningBudgetToAnchorFourHundred{}.

\begin{table}[htbp]
\caption{Search space of the tuning, identical for S, both injection routes of D and P1
(\texttt{configs/config.yaml}, \texttt{tuning.search\_space}).}
\label{tab:space}
\centering\small
\begin{tabular}{@{}lll@{}}
\toprule
Setting & Distribution & Range or choices \\
\midrule
Learning rate & log-uniform & \numCfgTuningSearchSpaceLearningRateLow{} to \numCfgTuningSearchSpaceLearningRateHigh{} \\
Weight decay & log-uniform & \numCfgTuningSearchSpaceWeightDecayLow{} to \numCfgTuningSearchSpaceWeightDecayHigh{} \\
Hidden width & categorical & \numCfgTuningSearchSpaceHiddenWidthChoicesItemOne{}, \numCfgTuningSearchSpaceHiddenWidthChoicesItemTwo{} \\
MEGNet blocks & categorical & \numCfgTuningSearchSpaceMegnetBlocksChoicesItemOne{}, \numCfgTuningSearchSpaceMegnetBlocksChoicesItemTwo{}, \numCfgTuningSearchSpaceMegnetBlocksChoicesItemThree{} \\
Dropout & uniform & \numCfgTuningSearchSpaceDropoutLow{} to \numCfgTuningSearchSpaceDropoutHigh{} \\
Batch size (sites) & categorical & \numCfgTuningSearchSpaceBatchSizeChoicesItemOne{}, \numCfgTuningSearchSpaceBatchSizeChoicesItemTwo{}, \numCfgTuningSearchSpaceBatchSizeChoicesItemThree{} \\
Readout width & categorical & \numCfgTuningSearchSpaceReadoutMlpWidthChoicesItemOne{}, \numCfgTuningSearchSpaceReadoutMlpWidthChoicesItemTwo{} \\
Pooling of the atoms & categorical & Set2Set, mean \\
\bottomrule
\end{tabular}
\end{table}

The originally registered network (v1) was tuned under the frozen plan, and its values are those of the tag
\tagtuned{} (Table~\ref{tab:tuned_v1}). The two injection routes of D were tuned separately, and the rule logged
before tuning kept the route with the lower mean, over the anchors, of the best validation error: D-state, with
\numTuningSummaryDecisionMeansEVDState{}~eV against \numTuningSummaryDecisionMeansEVDLate{}~eV for D-late, a
difference of \numTuningSummaryDecisionDifferenceEV{}~eV. The two routes are indistinguishable on this objective,
and the selection is not evidence that one route encodes the electronic structure better. The primary network
(v2) differs from v1 only in its initial weights (He-uniform initialisation and zero biases); it was adopted after
test-set diagnostics of v1 and is post hoc. S and D of v2 were retuned with the same protocol, and test hosts were
never loaded in its design screen or its tuning (Table~\ref{tab:tuned_v2}). The descriptor predictor of v2, P1-v2,
was not retuned: it uses the v1 values of P1 with the v2 initialisation. Across both tables, the settings chosen at
the smallest anchor rest on very few validation hosts and should be read as one reasonable setting rather than
an optimum.

\begin{table}[htbp]
\caption{Tuned values of the originally registered network (v1; \texttt{configs/tuned\_v1.yaml}, tag \tagtuned{}),
with the best validation MAE (eV for S and D; the mean absolute error over the standardised descriptors,
dimensionless, for P1) of each study (\texttt{results/tuning\_summary.json}). D is the selected D-state route;
D-late is reported for completeness.}
\label{tab:tuned_v1}
\centering\footnotesize\tighttabsb
\begin{adjustbox}{max width=\linewidth}
% Generated by paper/si/tools/make_si_tables.py from configs/tuned_v1.yaml, results/tuning_summary.json. Do not edit.
\begin{tabular}{@{}llllllllllr@{}}
\toprule
Model & Anchor & Learning rate & Weight decay & Width & Blocks & Dropout & Batch & Readout & Pooling & Val.\ MAE \\
\midrule
S & \numCfgTuningAnchorsItemOne{} & \numTunedVOneModelsSFiftyLearningRateSig{} & \numTunedVOneModelsSFiftyWeightDecaySig{} & \numTunedVOneModelsSFiftyHiddenWidth{} & \numTunedVOneModelsSFiftyMegnetBlocks{} & \numTunedVOneModelsSFiftyDropoutSig{} & \numTunedVOneModelsSFiftyBatchSize{} & \numTunedVOneModelsSFiftyReadoutMlpWidth{} & set2set & \numTuningSummaryStudiesSAFiftyBestValue{} \\
 & \numCfgTuningAnchorsItemTwo{} & \numTunedVOneModelsSTwoHundredLearningRateSig{} & \numTunedVOneModelsSTwoHundredWeightDecaySig{} & \numTunedVOneModelsSTwoHundredHiddenWidth{} & \numTunedVOneModelsSTwoHundredMegnetBlocks{} & \numTunedVOneModelsSTwoHundredDropoutSig{} & \numTunedVOneModelsSTwoHundredBatchSize{} & \numTunedVOneModelsSTwoHundredReadoutMlpWidth{} & set2set & \numTuningSummaryStudiesSATwoHundredBestValue{} \\
 & \numCfgTuningAnchorsItemThree{} & \numTunedVOneModelsSSixHundredFiftyFourLearningRateSig{} & \numTunedVOneModelsSSixHundredFiftyFourWeightDecaySig{} & \numTunedVOneModelsSSixHundredFiftyFourHiddenWidth{} & \numTunedVOneModelsSSixHundredFiftyFourMegnetBlocks{} & \numTunedVOneModelsSSixHundredFiftyFourDropoutSig{} & \numTunedVOneModelsSSixHundredFiftyFourBatchSize{} & \numTunedVOneModelsSSixHundredFiftyFourReadoutMlpWidth{} & mean & \numTuningSummaryStudiesSASixHundredFiftyFourBestValue{} \\
\addlinespace
D (D-state) & \numCfgTuningAnchorsItemOne{} & \numTunedVOneModelsDStateFiftyLearningRateSig{} & \numTunedVOneModelsDStateFiftyWeightDecaySig{} & \numTunedVOneModelsDStateFiftyHiddenWidth{} & \numTunedVOneModelsDStateFiftyMegnetBlocks{} & \numTunedVOneModelsDStateFiftyDropoutSig{} & \numTunedVOneModelsDStateFiftyBatchSize{} & \numTunedVOneModelsDStateFiftyReadoutMlpWidth{} & set2set & \numTuningSummaryStudiesDStateAFiftyBestValue{} \\
 & \numCfgTuningAnchorsItemTwo{} & \numTunedVOneModelsDStateTwoHundredLearningRateSig{} & \numTunedVOneModelsDStateTwoHundredWeightDecaySig{} & \numTunedVOneModelsDStateTwoHundredHiddenWidth{} & \numTunedVOneModelsDStateTwoHundredMegnetBlocks{} & \numTunedVOneModelsDStateTwoHundredDropoutSig{} & \numTunedVOneModelsDStateTwoHundredBatchSize{} & \numTunedVOneModelsDStateTwoHundredReadoutMlpWidth{} & set2set & \numTuningSummaryStudiesDStateATwoHundredBestValue{} \\
 & \numCfgTuningAnchorsItemThree{} & \numTunedVOneModelsDStateSixHundredFiftyFourLearningRateSig{} & \numTunedVOneModelsDStateSixHundredFiftyFourWeightDecaySig{} & \numTunedVOneModelsDStateSixHundredFiftyFourHiddenWidth{} & \numTunedVOneModelsDStateSixHundredFiftyFourMegnetBlocks{} & \numTunedVOneModelsDStateSixHundredFiftyFourDropoutSig{} & \numTunedVOneModelsDStateSixHundredFiftyFourBatchSize{} & \numTunedVOneModelsDStateSixHundredFiftyFourReadoutMlpWidth{} & mean & \numTuningSummaryStudiesDStateASixHundredFiftyFourBestValue{} \\
\addlinespace
D-late & \numCfgTuningAnchorsItemOne{} & \numTunedVOneModelsDLateFiftyLearningRateSig{} & \numTunedVOneModelsDLateFiftyWeightDecaySig{} & \numTunedVOneModelsDLateFiftyHiddenWidth{} & \numTunedVOneModelsDLateFiftyMegnetBlocks{} & \numTunedVOneModelsDLateFiftyDropoutSig{} & \numTunedVOneModelsDLateFiftyBatchSize{} & \numTunedVOneModelsDLateFiftyReadoutMlpWidth{} & mean & \numTuningSummaryStudiesDLateAFiftyBestValue{} \\
 & \numCfgTuningAnchorsItemTwo{} & \numTunedVOneModelsDLateTwoHundredLearningRateSig{} & \numTunedVOneModelsDLateTwoHundredWeightDecaySig{} & \numTunedVOneModelsDLateTwoHundredHiddenWidth{} & \numTunedVOneModelsDLateTwoHundredMegnetBlocks{} & \numTunedVOneModelsDLateTwoHundredDropoutSig{} & \numTunedVOneModelsDLateTwoHundredBatchSize{} & \numTunedVOneModelsDLateTwoHundredReadoutMlpWidth{} & set2set & \numTuningSummaryStudiesDLateATwoHundredBestValue{} \\
 & \numCfgTuningAnchorsItemThree{} & \numTunedVOneModelsDLateSixHundredFiftyFourLearningRateSig{} & \numTunedVOneModelsDLateSixHundredFiftyFourWeightDecaySig{} & \numTunedVOneModelsDLateSixHundredFiftyFourHiddenWidth{} & \numTunedVOneModelsDLateSixHundredFiftyFourMegnetBlocks{} & \numTunedVOneModelsDLateSixHundredFiftyFourDropoutSig{} & \numTunedVOneModelsDLateSixHundredFiftyFourBatchSize{} & \numTunedVOneModelsDLateSixHundredFiftyFourReadoutMlpWidth{} & mean & \numTuningSummaryStudiesDLateASixHundredFiftyFourBestValue{} \\
\addlinespace
P1 & \numCfgTuningAnchorsItemOne{} & \numTunedVOneModelsPOneFiftyLearningRateSig{} & \numTunedVOneModelsPOneFiftyWeightDecaySig{} & \numTunedVOneModelsPOneFiftyHiddenWidth{} & \numTunedVOneModelsPOneFiftyMegnetBlocks{} & \numTunedVOneModelsPOneFiftyDropoutSig{} & \numTunedVOneModelsPOneFiftyBatchSize{} & \numTunedVOneModelsPOneFiftyReadoutMlpWidth{} & set2set & \numTuningSummaryStudiesPOneAFiftyBestValue{} \\
 & \numCfgTuningAnchorsItemTwo{} & \numTunedVOneModelsPOneTwoHundredLearningRateSig{} & \numTunedVOneModelsPOneTwoHundredWeightDecaySig{} & \numTunedVOneModelsPOneTwoHundredHiddenWidth{} & \numTunedVOneModelsPOneTwoHundredMegnetBlocks{} & \numTunedVOneModelsPOneTwoHundredDropoutSig{} & \numTunedVOneModelsPOneTwoHundredBatchSize{} & \numTunedVOneModelsPOneTwoHundredReadoutMlpWidth{} & mean & \numTuningSummaryStudiesPOneATwoHundredBestValue{} \\
 & \numCfgTuningAnchorsItemThree{} & \numTunedVOneModelsPOneSixHundredFiftyFourLearningRateSig{} & \numTunedVOneModelsPOneSixHundredFiftyFourWeightDecaySig{} & \numTunedVOneModelsPOneSixHundredFiftyFourHiddenWidth{} & \numTunedVOneModelsPOneSixHundredFiftyFourMegnetBlocks{} & \numTunedVOneModelsPOneSixHundredFiftyFourDropoutSig{} & \numTunedVOneModelsPOneSixHundredFiftyFourBatchSize{} & \numTunedVOneModelsPOneSixHundredFiftyFourReadoutMlpWidth{} & mean & \numTuningSummaryStudiesPOneASixHundredFiftyFourBestValue{} \\
\bottomrule
\end{tabular}

\end{adjustbox}
\end{table}

\begin{table}[htbp]
\caption{Tuned values of the primary network (v2, post hoc; \texttt{configs/tuned\_v2.yaml}). Same protocol and
search space as Table~\ref{tab:tuned_v1}; P1-v2 reuses the P1 values of Table~\ref{tab:tuned_v1}.}
\label{tab:tuned_v2}
\centering\footnotesize\tighttabsb
\begin{adjustbox}{max width=\linewidth}
% Generated by paper/si/tools/make_si_tables.py from configs/tuned_v2.yaml. Do not edit.
\begin{tabular}{@{}llllllllll@{}}
\toprule
Model & Anchor & Learning rate & Weight decay & Width & Blocks & Dropout & Batch & Readout & Pooling \\
\midrule
S & \numCfgTuningAnchorsItemOne{} & \numTunedVTwoModelsSFiftyLearningRateSig{} & \numTunedVTwoModelsSFiftyWeightDecaySig{} & \numTunedVTwoModelsSFiftyHiddenWidth{} & \numTunedVTwoModelsSFiftyMegnetBlocks{} & \numTunedVTwoModelsSFiftyDropoutSig{} & \numTunedVTwoModelsSFiftyBatchSize{} & \numTunedVTwoModelsSFiftyReadoutMlpWidth{} & mean \\
 & \numCfgTuningAnchorsItemTwo{} & \numTunedVTwoModelsSTwoHundredLearningRateSig{} & \numTunedVTwoModelsSTwoHundredWeightDecaySig{} & \numTunedVTwoModelsSTwoHundredHiddenWidth{} & \numTunedVTwoModelsSTwoHundredMegnetBlocks{} & \numTunedVTwoModelsSTwoHundredDropoutSig{} & \numTunedVTwoModelsSTwoHundredBatchSize{} & \numTunedVTwoModelsSTwoHundredReadoutMlpWidth{} & mean \\
 & \numCfgTuningAnchorsItemThree{} & \numTunedVTwoModelsSSixHundredFiftyFourLearningRateSig{} & \numTunedVTwoModelsSSixHundredFiftyFourWeightDecaySig{} & \numTunedVTwoModelsSSixHundredFiftyFourHiddenWidth{} & \numTunedVTwoModelsSSixHundredFiftyFourMegnetBlocks{} & \numTunedVTwoModelsSSixHundredFiftyFourDropoutSig{} & \numTunedVTwoModelsSSixHundredFiftyFourBatchSize{} & \numTunedVTwoModelsSSixHundredFiftyFourReadoutMlpWidth{} & set2set \\
\addlinespace
D (D-state) & \numCfgTuningAnchorsItemOne{} & \numTunedVTwoModelsDStateFiftyLearningRateSig{} & \numTunedVTwoModelsDStateFiftyWeightDecaySig{} & \numTunedVTwoModelsDStateFiftyHiddenWidth{} & \numTunedVTwoModelsDStateFiftyMegnetBlocks{} & \numTunedVTwoModelsDStateFiftyDropoutSig{} & \numTunedVTwoModelsDStateFiftyBatchSize{} & \numTunedVTwoModelsDStateFiftyReadoutMlpWidth{} & mean \\
 & \numCfgTuningAnchorsItemTwo{} & \numTunedVTwoModelsDStateTwoHundredLearningRateSig{} & \numTunedVTwoModelsDStateTwoHundredWeightDecaySig{} & \numTunedVTwoModelsDStateTwoHundredHiddenWidth{} & \numTunedVTwoModelsDStateTwoHundredMegnetBlocks{} & \numTunedVTwoModelsDStateTwoHundredDropoutSig{} & \numTunedVTwoModelsDStateTwoHundredBatchSize{} & \numTunedVTwoModelsDStateTwoHundredReadoutMlpWidth{} & mean \\
 & \numCfgTuningAnchorsItemThree{} & \numTunedVTwoModelsDStateSixHundredFiftyFourLearningRateSig{} & \numTunedVTwoModelsDStateSixHundredFiftyFourWeightDecaySig{} & \numTunedVTwoModelsDStateSixHundredFiftyFourHiddenWidth{} & \numTunedVTwoModelsDStateSixHundredFiftyFourMegnetBlocks{} & \numTunedVTwoModelsDStateSixHundredFiftyFourDropoutSig{} & \numTunedVTwoModelsDStateSixHundredFiftyFourBatchSize{} & \numTunedVTwoModelsDStateSixHundredFiftyFourReadoutMlpWidth{} & mean \\
\bottomrule
\end{tabular}

\end{adjustbox}
\end{table}

\subsection{The cap on training length and its measured effects}

The cap of \numCfgTrainingMaxEpochs{} passes over the training hosts and the patience of
\numCfgTrainingEarlyStoppingPatience{} passes were set by a rule logged before the calibration pilot ran: the cap
is a round number of passes above the slowest convergence seen in the pilot, with a safety factor, and the
patience follows from the cap (\texttt{docs/deviations.md}, the \sessioncap{} rows; \texttt{docs/methods\_notes.md}). The
registered check of the cap retrained S and D with a cap of \numGOneCapSensitivityCapsItemTwo{} passes at
\numGOneCapSensitivityBudget{} hosts on \numGOneCapSensitivityBootstrapNResamples{} resamples and all seeds. All
\numGOneDiagnosticsEpochCapPairsNPairs{} pairs stopped at the same pass and gave identical predictions
(\numGOneDiagnosticsEpochCapPairsNIdentical{} of \numGOneDiagnosticsEpochCapPairsNPairs{}), so the advantage of the
descriptors was \numGOneCapSensitivityATwoHundredMean{}~eV under both caps
\ci{\numGOneCapSensitivityATwoHundredCiLo{}}{\numGOneCapSensitivityATwoHundredCiHi{}} and the primary result was
cap-insensitive by the registered criterion (registered v1; \texttt{results/g1\_cap\_sensitivity.json}).

The extensions show that the cap binds more often than those runs suggested, and mostly for the structure-only
network (Table~\ref{tab:capped}). Among the precision runs of v1 at \numCfgBudgetsHostsItemSix{} hosts,
\numPrecisionABFourHundredBSixHundredFiftyFourDescriptiveCappedPrecisionSixHundredFiftyFourVOneSNCapped{} of
\numPrecisionABFourHundredBSixHundredFiftyFourDescriptiveCappedPrecisionSixHundredFiftyFourVOneSNRuns{} S runs
reached the cap, against
\numPrecisionABFourHundredBSixHundredFiftyFourDescriptiveCappedPrecisionSixHundredFiftyFourVOneDNCapped{} D runs;
for v2 the counts were
\numPrecisionABFourHundredBSixHundredFiftyFourDescriptiveCappedPrecisionSixHundredFiftyFourVTwoSNCapped{} and
\numPrecisionABFourHundredBSixHundredFiftyFourDescriptiveCappedPrecisionSixHundredFiftyFourVTwoDNCapped{}. In every
group of S runs with capped members, the capped runs had the lower mean error, so the structure-only network may
still have been improving when it was stopped at the largest training sets. Because capped and uncapped runs differ
in resample and fold as well, Table~\ref{tab:capped} is descriptive and carries no test. The post hoc rerun of every
capped primary run with a longer cap is reported in Section~\ref{si:sensitivities}: it changes the advantage at
\numCfgBudgetsHostsItemSix{} hosts by
\numSensCappedPostHocPerBudgetSixHundredFiftyFourChangeTwinMinusPrimaryAMean{}~eV on v1, and no such reruns exist
for v2.

\begin{table}[htbp]
\caption{Runs that stopped at the cap of \numCfgTrainingMaxEpochs{} passes (extension, descriptive, no test), and
the mean single-run site MAE (eV) of capped and uncapped runs. Precision: the pooled runs at the two largest
budgets; grouped CV: all folds at each $K$. Records: \recprecision{} and \reccv{}, key
\texttt{payload.descriptive.capped}.}
\label{tab:capped}
\centering\scriptsize
% Generated by paper/si/tools/make_si_tables.py from results/precision_a_b400_b654.json, results/cv_extended_budget.json (payload.descriptive.capped). Do not edit.
\begin{tabular}{@{}llllrrrr@{}}
\toprule
Design & Budget or folds & Network & Model & Runs & Stopped at cap & MAE, capped & MAE, not capped \\
\midrule
precision & \numCfgBudgetsHostsItemFive{} & v2 & S & \numPrecisionABFourHundredBSixHundredFiftyFourDescriptiveCappedPrecisionFourHundredVTwoSNRuns{} & \numPrecisionABFourHundredBSixHundredFiftyFourDescriptiveCappedPrecisionFourHundredVTwoSNCapped{} & \numPrecisionABFourHundredBSixHundredFiftyFourDescriptiveCappedPrecisionFourHundredVTwoSMeanSingleRunMaeCapped{} & \numPrecisionABFourHundredBSixHundredFiftyFourDescriptiveCappedPrecisionFourHundredVTwoSMeanSingleRunMaeUncapped{} \\
precision &  & v2 & D & \numPrecisionABFourHundredBSixHundredFiftyFourDescriptiveCappedPrecisionFourHundredVTwoDNRuns{} & \numPrecisionABFourHundredBSixHundredFiftyFourDescriptiveCappedPrecisionFourHundredVTwoDNCapped{} & -- & \numPrecisionABFourHundredBSixHundredFiftyFourDescriptiveCappedPrecisionFourHundredVTwoDMeanSingleRunMaeUncapped{} \\
precision &  & v1 & S & \numPrecisionABFourHundredBSixHundredFiftyFourDescriptiveCappedPrecisionFourHundredVOneSNRuns{} & \numPrecisionABFourHundredBSixHundredFiftyFourDescriptiveCappedPrecisionFourHundredVOneSNCapped{} & \numPrecisionABFourHundredBSixHundredFiftyFourDescriptiveCappedPrecisionFourHundredVOneSMeanSingleRunMaeCapped{} & \numPrecisionABFourHundredBSixHundredFiftyFourDescriptiveCappedPrecisionFourHundredVOneSMeanSingleRunMaeUncapped{} \\
precision &  & v1 & D & \numPrecisionABFourHundredBSixHundredFiftyFourDescriptiveCappedPrecisionFourHundredVOneDNRuns{} & \numPrecisionABFourHundredBSixHundredFiftyFourDescriptiveCappedPrecisionFourHundredVOneDNCapped{} & \numPrecisionABFourHundredBSixHundredFiftyFourDescriptiveCappedPrecisionFourHundredVOneDMeanSingleRunMaeCapped{} & \numPrecisionABFourHundredBSixHundredFiftyFourDescriptiveCappedPrecisionFourHundredVOneDMeanSingleRunMaeUncapped{} \\
\addlinespace
precision & \numCfgBudgetsHostsItemSix{} & v2 & S & \numPrecisionABFourHundredBSixHundredFiftyFourDescriptiveCappedPrecisionSixHundredFiftyFourVTwoSNRuns{} & \numPrecisionABFourHundredBSixHundredFiftyFourDescriptiveCappedPrecisionSixHundredFiftyFourVTwoSNCapped{} & \numPrecisionABFourHundredBSixHundredFiftyFourDescriptiveCappedPrecisionSixHundredFiftyFourVTwoSMeanSingleRunMaeCapped{} & \numPrecisionABFourHundredBSixHundredFiftyFourDescriptiveCappedPrecisionSixHundredFiftyFourVTwoSMeanSingleRunMaeUncapped{} \\
precision &  & v2 & D & \numPrecisionABFourHundredBSixHundredFiftyFourDescriptiveCappedPrecisionSixHundredFiftyFourVTwoDNRuns{} & \numPrecisionABFourHundredBSixHundredFiftyFourDescriptiveCappedPrecisionSixHundredFiftyFourVTwoDNCapped{} & \numPrecisionABFourHundredBSixHundredFiftyFourDescriptiveCappedPrecisionSixHundredFiftyFourVTwoDMeanSingleRunMaeCapped{} & \numPrecisionABFourHundredBSixHundredFiftyFourDescriptiveCappedPrecisionSixHundredFiftyFourVTwoDMeanSingleRunMaeUncapped{} \\
precision &  & v1 & S & \numPrecisionABFourHundredBSixHundredFiftyFourDescriptiveCappedPrecisionSixHundredFiftyFourVOneSNRuns{} & \numPrecisionABFourHundredBSixHundredFiftyFourDescriptiveCappedPrecisionSixHundredFiftyFourVOneSNCapped{} & \numPrecisionABFourHundredBSixHundredFiftyFourDescriptiveCappedPrecisionSixHundredFiftyFourVOneSMeanSingleRunMaeCapped{} & \numPrecisionABFourHundredBSixHundredFiftyFourDescriptiveCappedPrecisionSixHundredFiftyFourVOneSMeanSingleRunMaeUncapped{} \\
precision &  & v1 & D & \numPrecisionABFourHundredBSixHundredFiftyFourDescriptiveCappedPrecisionSixHundredFiftyFourVOneDNRuns{} & \numPrecisionABFourHundredBSixHundredFiftyFourDescriptiveCappedPrecisionSixHundredFiftyFourVOneDNCapped{} & \numPrecisionABFourHundredBSixHundredFiftyFourDescriptiveCappedPrecisionSixHundredFiftyFourVOneDMeanSingleRunMaeCapped{} & \numPrecisionABFourHundredBSixHundredFiftyFourDescriptiveCappedPrecisionSixHundredFiftyFourVOneDMeanSingleRunMaeUncapped{} \\
\addlinespace
grouped CV & $K = \numCvExtendedBudgetDescriptiveCappedCvTenVTwoSPoint{}$ & v2 & S & \numCvExtendedBudgetDescriptiveCappedCvTenVTwoSNRuns{} & \numCvExtendedBudgetDescriptiveCappedCvTenVTwoSNCapped{} & \numCvExtendedBudgetDescriptiveCappedCvTenVTwoSMeanSingleRunMaeCapped{} & \numCvExtendedBudgetDescriptiveCappedCvTenVTwoSMeanSingleRunMaeUncapped{} \\
grouped CV &  & v2 & D & \numCvExtendedBudgetDescriptiveCappedCvTenVTwoDNRuns{} & \numCvExtendedBudgetDescriptiveCappedCvTenVTwoDNCapped{} & \numCvExtendedBudgetDescriptiveCappedCvTenVTwoDMeanSingleRunMaeCapped{} & \numCvExtendedBudgetDescriptiveCappedCvTenVTwoDMeanSingleRunMaeUncapped{} \\
grouped CV &  & v1 & S & \numCvExtendedBudgetDescriptiveCappedCvTenVOneSNRuns{} & \numCvExtendedBudgetDescriptiveCappedCvTenVOneSNCapped{} & \numCvExtendedBudgetDescriptiveCappedCvTenVOneSMeanSingleRunMaeCapped{} & \numCvExtendedBudgetDescriptiveCappedCvTenVOneSMeanSingleRunMaeUncapped{} \\
grouped CV &  & v1 & D & \numCvExtendedBudgetDescriptiveCappedCvTenVOneDNRuns{} & \numCvExtendedBudgetDescriptiveCappedCvTenVOneDNCapped{} & \numCvExtendedBudgetDescriptiveCappedCvTenVOneDMeanSingleRunMaeCapped{} & \numCvExtendedBudgetDescriptiveCappedCvTenVOneDMeanSingleRunMaeUncapped{} \\
\addlinespace
grouped CV & $K = \numCvExtendedBudgetDescriptiveCappedCvTwentyVTwoSPoint{}$ & v2 & S & \numCvExtendedBudgetDescriptiveCappedCvTwentyVTwoSNRuns{} & \numCvExtendedBudgetDescriptiveCappedCvTwentyVTwoSNCapped{} & \numCvExtendedBudgetDescriptiveCappedCvTwentyVTwoSMeanSingleRunMaeCapped{} & \numCvExtendedBudgetDescriptiveCappedCvTwentyVTwoSMeanSingleRunMaeUncapped{} \\
grouped CV &  & v2 & D & \numCvExtendedBudgetDescriptiveCappedCvTwentyVTwoDNRuns{} & \numCvExtendedBudgetDescriptiveCappedCvTwentyVTwoDNCapped{} & \numCvExtendedBudgetDescriptiveCappedCvTwentyVTwoDMeanSingleRunMaeCapped{} & \numCvExtendedBudgetDescriptiveCappedCvTwentyVTwoDMeanSingleRunMaeUncapped{} \\
grouped CV &  & v1 & S & \numCvExtendedBudgetDescriptiveCappedCvTwentyVOneSNRuns{} & \numCvExtendedBudgetDescriptiveCappedCvTwentyVOneSNCapped{} & \numCvExtendedBudgetDescriptiveCappedCvTwentyVOneSMeanSingleRunMaeCapped{} & \numCvExtendedBudgetDescriptiveCappedCvTwentyVOneSMeanSingleRunMaeUncapped{} \\
grouped CV &  & v1 & D & \numCvExtendedBudgetDescriptiveCappedCvTwentyVOneDNRuns{} & \numCvExtendedBudgetDescriptiveCappedCvTwentyVOneDNCapped{} & \numCvExtendedBudgetDescriptiveCappedCvTwentyVOneDMeanSingleRunMaeCapped{} & \numCvExtendedBudgetDescriptiveCappedCvTwentyVOneDMeanSingleRunMaeUncapped{} \\
\bottomrule
\end{tabular}

\end{table}

\subsection{Hardware}

The registered tuning and the registered sweep of the originally registered network (v1) ran on one
\gpulfortys{}, with several training processes sharing the card; a process started only when the memory of the
running jobs left a fixed margin (\texttt{docs/methods\_notes.md}). These settings change wall time only, not any
setting, split, seed or stopping rule. The run records of the original sweep, and those of the collaborator's v2
tuning and production, name a CUDA device but not the GPU model. The extension training runs (precision, grouped
cross-validation, high-moment, CGCNN and capped twins) ran on the GPUs of the National Research Platform, one
training process per GPU, with PyTorch \vertorch{} and matgl \vermatgl{}; their records name the GPU model:
\gpuaten{}, \gputhreeoh{}, \gpulfour{} and \gpufourty{}. The inference stages of the latent probe and the geometry
evaluation also used an \gputitan{} (\texttt{docs/nrp/n3\_ingest.md}). Training on a GPU is not bitwise
reproducible across GPU models, because the summation order of the aggregation over neighbours differs;
Section~\ref{si:hardware} gives the repeatability study and the size of the effect on the predicted formation
energies.

\section{Split designs}
\label{si:splits}

Every split keeps all oxygen sites of one host oxide on the same side, so no chemically identical compound is
seen in training and in testing; in this database one formula is one host, so grouping by formula and by host
coincide. The data that enter the splits are those of the filter cascade (Figure~\ref{fig:cascade}; counts in
Table~\ref{tab:cascade}), and the four designs are drawn in Figure~\ref{fig:splits}.

\begin{figure}[htbp]
\centering
\sifig{fig_cascade}
\caption{Filter cascade that defines universe v1. From the completed neutral oxygen-vacancy entries of the
release of Kumagai et al., the cascade keeps their neutral machine-learning subset, which already excludes the
perturbed-host-state entries, the entries with undetermined band edges and the dynamically unstable hosts. Three
filters follow in order: D1 removes the entries flagged as a split vacancy, an unknown defect type or a
configuration different from the initial one; D2 removes the host Sr$_2$Zr$_7$O$_{16}$, whose site labels do not
match the targets of the release; D3 keeps an entry of the two hosts with inconsistent cell metadata only if its
vacancy site maps to an oxygen atom with a matching first cation shell, and removes nothing. Dashed boxes are the
entries excluded at each step. Site and host counts after every step are in Table~\ref{tab:cascade}.}
\label{fig:cascade}
\end{figure}

In the primary protocol each of \numCfgSplitNOuterResamples{} outer resamples holds out
\numSplitMetaNTestHosts{} test hosts, a fraction \numSplitMetaTestFraction{} of the universe, stratified by the
host-mean formation energy, and orders the remaining pool of \numSplitMetaPoolSize{} hosts so that the training
set at budget $B$ is its first $B$ hosts. The budgets, \numCfgBudgetsHostsItemOne{} to \numCfgBudgetsHostsItemSix{}
hosts, are therefore nested within a resample, and the test hosts of a resample are the same at every budget, so
that the learning curves compare budgets on identical test oxides. Across the resamples,
\numSplitMetaCoverageInTestAtLeastOnce{} hosts are tested at least once and \numSplitMetaCoverageNeverInTest{}
never. The precision extension adds resamples drawn by the same rule, to \numPrecisionABFourHundredBSixHundredFiftyFourSubsetsThirtyXFiveNResamples{}
resamples with \numPrecisionABFourHundredBSixHundredFiftyFourSubsetsThirtyXFiveNSeeds{} training seeds each at
the two largest budgets (extension).

The published split of Kiyohara et al.\ \citep{Kiyohara2025machinelearning} is used once, to compare the
structure-only network with the published accuracy: the networks train on its
\numFilterCascadeKiyoharaSplitHostsTrain{} training hosts, stop early on its \numFilterCascadeKiyoharaSplitHostsVal{}
validation hosts and are tested on its \numFilterCascadeKiyoharaSplitHostsTest{} test hosts
(\numFilterCascadeKiyoharaSplitSitesTest{} sites). This comparison is not fully held out at the level of the
settings. The tuning used the nested training sets of resample \numCfgTuningAnchorResample{}, which was drawn
independently of the published split, so \numCZeroOfficialTuningOverlapWithKiyoharaTestSixHundredFiftyFourInKiyoharaTest{}
of the published test hosts lie in the largest tuning anchor
(\numCZeroOfficialTuningOverlapWithKiyoharaTestSixHundredFiftyFourOfWhichValidationHosts{} of them among its
validation hosts), and the comparison uses the settings of that anchor. The network weights never saw the labels
of those hosts in the comparison itself (\texttt{results/c0\_official.json}, key
\texttt{tuning\_overlap\_with\_kiyohara\_test}).

The grouped $K$-fold cross-validation (extension) partitions all \numSplitCvNHosts{} hosts into
$K = \numSplitCvKsItemOne{}$ or $K = \numSplitCvKsItemTwo{}$ folds, so that every host is tested exactly once and the
mean training set, \numCvGOneVOneXTrainingHostsKTen{} and \numCvGOneVOneXTrainingHostsKTwenty{} hosts, exceeds the
largest registered budget. Its test population is the whole universe, including the
\numSplitMetaCoverageNeverInTest{} hosts that no outer resample tests, so a cross-validated error is not directly
comparable with a registered budget at a similar number of training hosts.

Leave-one-chemistry-out asks whether the models transfer to chemistries absent from training. Each host belongs to
the family given by the set of periodic-table groups of its cations, which gives \numSplitLocoNFamilies{} families
over the universe, and \numSplitLocoNFolds{} folds each hold out whole families. The runs of the originally
registered network follow the registered design but are the collaborator's, whose validation carve-out differs
from the one our later clarification specifies; they are adopted as run and labelled extension at the results
freeze (\texttt{docs/freeze\_notes.md}), while \texttt{docs/provenance\_labels.md} records them as registered in
design and leaves the choice to the author. The runs of the primary network are post hoc.

\begin{figure}[htbp]
\centering
\sifig{fig_splits}
\caption{Split designs, drawn as sets of hosts (dark: test; light: training; mid-grey: validation). Every split
keeps all sites of a host on one side. Widths are schematic, not to scale. (\textbf{a}) Primary protocol: each
outer resample holds out \numSplitMetaNTestHosts{} test hosts, stratified by host-mean $E_\mathrm{f}$, and orders
the remaining pool so that the training set at budget $B$ is the first $B$ hosts; budgets are nested, and the
learning curves are drawn against $B$. Validation hosts are carved out of each training set. (\textbf{b}) The
published split of Kiyohara et al.: the networks train on its training hosts, stop early on its validation hosts
and are tested on its test hosts. (\textbf{c}) Grouped $K$-fold cross-validation ($K = \numSplitCvKsItemOne{}$ and
$\numSplitCvKsItemTwo{}$; extension): every host is tested exactly once, with training sets larger than the largest
budget of (a). (\textbf{d}) Leave-one-chemistry-out: hosts are grouped into chemistry families (the set of
periodic-table groups of the cations), and \numSplitLocoNFolds{} folds, each holding out whole families, test
transfer to chemistries absent from training.}
\label{fig:splits}
\end{figure}
\FloatBarrier

\section{Site resolution within a host, by chemistry and site}
\label{si:withinhost}

This section asks whether the ability of the networks to order the oxygen sites of one host depends on the
chemistry of the host or on the site (extension, not registered; both networks at \numWithinHostBreakdownBudget{}
training hosts). The rules (groups, thresholds, metrics and a host-cluster bootstrap paired across models and networks) were logged in \texttt{docs/deviations.md}
before any number was computed, and the analysis reads the existing test predictions without training anything
(\texttt{results/within\_host\_breakdown.json}; \texttt{docs/within\_host\_physics.md}). The measure is the
within-host residual error: for a host with at least two tested sites, the prediction and the DFT energy of each
site are taken relative to their host means, and the mean absolute difference of these relative values is the
error that remains once every host-level offset is removed. Its reference is the zero-skill value, the mean
absolute deviation of the DFT energies from their host mean, which is the within-host error of any predictor that
is constant within a host, such as the physics floor B0 or a network blind to the site. A network has
within-host skill in a group when its within-host error lies below the zero-skill value.

Each site's prediction is the seed ensemble averaged over the resamples in which its host was tested, so the
breakdown covers \numWithinHostBreakdownGroupsPooledAllNHosts{} hosts and \numWithinHostBreakdownGroupsPooledAllNSites{}
sites. This averaging differs from the registered estimator, which averages errors per resample: it lowers the
pooled site MAE of the originally registered network by
\numWithinHostBreakdownSanityISiteMaeVOneSiteAveragedMinusRecordS{}~eV for S and
\numWithinHostBreakdownSanityISiteMaeVOneSiteAveragedMinusRecordD{}~eV for D, so the values of this section are
comparable with each other but not with the registered learning curves.

Over all hosts, both networks resolve sites within a host, with within-host errors of
\numWithinHostBreakdownGroupsPooledAllVTwoSWithinHostMaePoint{}~eV (S) and
\numWithinHostBreakdownGroupsPooledAllVTwoDWithinHostMaePoint{}~eV (D) on the primary network (v2, post hoc)
against a zero-skill value of \numWithinHostBreakdownGroupsPooledAllVTwoSZeroSkillMaePoint{}~eV
(Table~\ref{tab:withinhost}). On v2 the descriptors leave the within-host error unchanged, S $-$ D =
\numWithinHostBreakdownGroupsPooledAllDifferencesSMinusDVTwoWithinHostMaePoint{}~eV
\ci{\numWithinHostBreakdownGroupsPooledAllDifferencesSMinusDVTwoWithinHostMaeCiNinetyFiveItemOne{}}{\numWithinHostBreakdownGroupsPooledAllDifferencesSMinusDVTwoWithinHostMaeCiNinetyFiveItemTwo{}};
on the originally registered network (v1) D has the larger within-host error, S $-$ D =
\numWithinHostBreakdownGroupsPooledAllDifferencesSMinusDVOneWithinHostMaePoint{}~eV
\ci{\numWithinHostBreakdownGroupsPooledAllDifferencesSMinusDVOneWithinHostMaeCiNinetyFiveItemOne{}}{\numWithinHostBreakdownGroupsPooledAllDifferencesSMinusDVOneWithinHostMaeCiNinetyFiveItemTwo{}},
so on v1 the network given the descriptors orders the sites of a host less well than the structure-only network.
The chemistry enters mainly through the spread of the target itself. Hosts whose cations have an empty $d$ shell
(\dzero{}) have a smaller site-to-site spread than hosts with a filled $d$ shell (\dten{}), with zero-skill values of
\numWithinHostBreakdownGroupsDClassDZeroVTwoSZeroSkillMaePoint{} and
\numWithinHostBreakdownGroupsDClassDTenVTwoSZeroSkillMaePoint{}~eV, and both networks keep within-host skill in both
classes. The v1 penalty of the descriptors is concentrated in the \dten{} hosts
(S $-$ D = \numWithinHostBreakdownGroupsDClassDTenDifferencesSMinusDVOneWithinHostMaePoint{}~eV
\ci{\numWithinHostBreakdownGroupsDClassDTenDifferencesSMinusDVOneWithinHostMaeCiNinetyFiveItemOne{}}{\numWithinHostBreakdownGroupsDClassDTenDifferencesSMinusDVOneWithinHostMaeCiNinetyFiveItemTwo{}})
and in the hosts that contain a cation with a high perturbed-host-state rate (Nb, W, Ti, V, Ta or Mo;
\numWithinHostBreakdownGroupsPhsPronePhsProneDifferencesSMinusDVOneWithinHostMaePoint{}~eV
\ci{\numWithinHostBreakdownGroupsPhsPronePhsProneDifferencesSMinusDVOneWithinHostMaeCiNinetyFiveItemOne{}}{\numWithinHostBreakdownGroupsPhsPronePhsProneDifferencesSMinusDVOneWithinHostMaeCiNinetyFiveItemTwo{}}).
Across the thirds of host stability the within-host error barely changes, while the site error of S on v2 rises
from \numWithinHostBreakdownGroupsStabilityTercileLowerVTwoSMaePoint{}~eV in the most stable third to
\numWithinHostBreakdownGroupsStabilityTercileUpperVTwoSMaePoint{}~eV in the least stable third. This pattern is
consistent with an error in the host-level offset rather than in the ordering of sites, a reading that is not
tested here. Among the thirds of the band gap, the smallest-gap third has the largest within-host errors of both
networks, and the errors at the most and least negative Bader charges of the vacancy oxygen within a host are
close to each other (Figure~\ref{fig:si_withinhost_groups}).

These are descriptive comparisons. The breakdown reports an interval for every group, model, network and metric
without correction for multiplicity, so some intervals exclude zero by chance alone, and an isolated interval that
excludes zero is weak evidence; groups with fewer hosts than the small-group threshold carry no interval. The
cation groups overlap, because a host with several cations is counted under each, and they are confounded with
each other and with structure type. Every group summary describes the filtered universe, that is, deep or
localised neutral vacancies only. Physical readings of these patterns are hypotheses
(\texttt{docs/within\_host\_physics.md}) and are not tested here.

\begin{table}[htbp]
\caption{Within-host residual MAE (eV) at \numWithinHostBreakdownBudget{} training hosts by group (extension;
both networks; v2 post hoc). Hosts and sites: those with at least two tested sites, on which the within-host error
is defined. Zero skill: the within-host error of any predictor constant within a host. Point values; the
two-sided intervals of the host-cluster bootstrap, and the per-cation groups of Figure~\ref{fig:si_withinhost_groups}a, are
in \texttt{results/within\_host\_breakdown.json} (\texttt{payload.groups}). The unassigned $d$ class is below the
small-group threshold and is descriptive only.}
\label{tab:withinhost}
\centering\small
% Generated by paper/si/tools/make_si_tables.py from results/within_host_breakdown.json (payload.groups). Do not edit.
\begin{tabular}{@{}lrrrrrrr@{}}
\toprule
 & & & & \pairhead{Primary (v2)} & \pairheadlast{Registered (v1)} \\
\whrules
Group & Hosts & Sites & Zero skill & S & D & S & D \\
\midrule
All tested hosts & \numWithinHostBreakdownGroupsPooledAllNHostsByMetricWithinHostMae{} & \numWithinHostBreakdownGroupsPooledAllNSitesWithin{} & \numWithinHostBreakdownGroupsPooledAllVTwoSZeroSkillMaePoint{} & \numWithinHostBreakdownGroupsPooledAllVTwoSWithinHostMaePoint{} & \numWithinHostBreakdownGroupsPooledAllVTwoDWithinHostMaePoint{} & \numWithinHostBreakdownGroupsPooledAllVOneSWithinHostMaePoint{} & \numWithinHostBreakdownGroupsPooledAllVOneDWithinHostMaePoint{} \\
\addlinespace
\whgroup{(b) $d$-electron class of the host} \\
$d$ class \dzero{} & \numWithinHostBreakdownGroupsDClassDZeroNHostsByMetricWithinHostMae{} & \numWithinHostBreakdownGroupsDClassDZeroNSitesWithin{} & \numWithinHostBreakdownGroupsDClassDZeroVTwoSZeroSkillMaePoint{} & \numWithinHostBreakdownGroupsDClassDZeroVTwoSWithinHostMaePoint{} & \numWithinHostBreakdownGroupsDClassDZeroVTwoDWithinHostMaePoint{} & \numWithinHostBreakdownGroupsDClassDZeroVOneSWithinHostMaePoint{} & \numWithinHostBreakdownGroupsDClassDZeroVOneDWithinHostMaePoint{} \\
$d$ class \dten{} & \numWithinHostBreakdownGroupsDClassDTenNHostsByMetricWithinHostMae{} & \numWithinHostBreakdownGroupsDClassDTenNSitesWithin{} & \numWithinHostBreakdownGroupsDClassDTenVTwoSZeroSkillMaePoint{} & \numWithinHostBreakdownGroupsDClassDTenVTwoSWithinHostMaePoint{} & \numWithinHostBreakdownGroupsDClassDTenVTwoDWithinHostMaePoint{} & \numWithinHostBreakdownGroupsDClassDTenVOneSWithinHostMaePoint{} & \numWithinHostBreakdownGroupsDClassDTenVOneDWithinHostMaePoint{} \\
$d$ class, other & \numWithinHostBreakdownGroupsDClassOtherNHostsByMetricWithinHostMae{} & \numWithinHostBreakdownGroupsDClassOtherNSitesWithin{} & \numWithinHostBreakdownGroupsDClassOtherVTwoSZeroSkillMaePoint{} & \numWithinHostBreakdownGroupsDClassOtherVTwoSWithinHostMaePoint{} & \numWithinHostBreakdownGroupsDClassOtherVTwoDWithinHostMaePoint{} & \numWithinHostBreakdownGroupsDClassOtherVOneSWithinHostMaePoint{} & \numWithinHostBreakdownGroupsDClassOtherVOneDWithinHostMaePoint{} \\
$d$ class, unassigned & \numWithinHostBreakdownGroupsDClassUnassignedNHostsByMetricWithinHostMae{} & \numWithinHostBreakdownGroupsDClassUnassignedNSitesWithin{} & \numWithinHostBreakdownGroupsDClassUnassignedVTwoSZeroSkillMaePoint{} & \numWithinHostBreakdownGroupsDClassUnassignedVTwoSWithinHostMaePoint{} & \numWithinHostBreakdownGroupsDClassUnassignedVTwoDWithinHostMaePoint{} & \numWithinHostBreakdownGroupsDClassUnassignedVOneSWithinHostMaePoint{} & \numWithinHostBreakdownGroupsDClassUnassignedVOneDWithinHostMaePoint{} \\
\addlinespace
\whgroup{(c) cations with high perturbed-host-state rates} \\
\begin{tabular}[t]{@{}l@{}}perturbed-host-state-prone\\cation present\end{tabular} & \numWithinHostBreakdownGroupsPhsPronePhsProneNHostsByMetricWithinHostMae{} & \numWithinHostBreakdownGroupsPhsPronePhsProneNSitesWithin{} & \numWithinHostBreakdownGroupsPhsPronePhsProneVTwoSZeroSkillMaePoint{} & \numWithinHostBreakdownGroupsPhsPronePhsProneVTwoSWithinHostMaePoint{} & \numWithinHostBreakdownGroupsPhsPronePhsProneVTwoDWithinHostMaePoint{} & \numWithinHostBreakdownGroupsPhsPronePhsProneVOneSWithinHostMaePoint{} & \numWithinHostBreakdownGroupsPhsPronePhsProneVOneDWithinHostMaePoint{} \\
\begin{tabular}[t]{@{}l@{}}no perturbed-host-state-\\prone cation\end{tabular} & \numWithinHostBreakdownGroupsPhsProneNotPhsProneNHostsByMetricWithinHostMae{} & \numWithinHostBreakdownGroupsPhsProneNotPhsProneNSitesWithin{} & \numWithinHostBreakdownGroupsPhsProneNotPhsProneVTwoSZeroSkillMaePoint{} & \numWithinHostBreakdownGroupsPhsProneNotPhsProneVTwoSWithinHostMaePoint{} & \numWithinHostBreakdownGroupsPhsProneNotPhsProneVTwoDWithinHostMaePoint{} & \numWithinHostBreakdownGroupsPhsProneNotPhsProneVOneSWithinHostMaePoint{} & \numWithinHostBreakdownGroupsPhsProneNotPhsProneVOneDWithinHostMaePoint{} \\
\addlinespace
\whgroup{(d) host stability (formation energy per atom)} \\
most stable third & \numWithinHostBreakdownGroupsStabilityTercileLowerNHostsByMetricWithinHostMae{} & \numWithinHostBreakdownGroupsStabilityTercileLowerNSitesWithin{} & \numWithinHostBreakdownGroupsStabilityTercileLowerVTwoSZeroSkillMaePoint{} & \numWithinHostBreakdownGroupsStabilityTercileLowerVTwoSWithinHostMaePoint{} & \numWithinHostBreakdownGroupsStabilityTercileLowerVTwoDWithinHostMaePoint{} & \numWithinHostBreakdownGroupsStabilityTercileLowerVOneSWithinHostMaePoint{} & \numWithinHostBreakdownGroupsStabilityTercileLowerVOneDWithinHostMaePoint{} \\
middle third & \numWithinHostBreakdownGroupsStabilityTercileMiddleNHostsByMetricWithinHostMae{} & \numWithinHostBreakdownGroupsStabilityTercileMiddleNSitesWithin{} & \numWithinHostBreakdownGroupsStabilityTercileMiddleVTwoSZeroSkillMaePoint{} & \numWithinHostBreakdownGroupsStabilityTercileMiddleVTwoSWithinHostMaePoint{} & \numWithinHostBreakdownGroupsStabilityTercileMiddleVTwoDWithinHostMaePoint{} & \numWithinHostBreakdownGroupsStabilityTercileMiddleVOneSWithinHostMaePoint{} & \numWithinHostBreakdownGroupsStabilityTercileMiddleVOneDWithinHostMaePoint{} \\
least stable third & \numWithinHostBreakdownGroupsStabilityTercileUpperNHostsByMetricWithinHostMae{} & \numWithinHostBreakdownGroupsStabilityTercileUpperNSitesWithin{} & \numWithinHostBreakdownGroupsStabilityTercileUpperVTwoSZeroSkillMaePoint{} & \numWithinHostBreakdownGroupsStabilityTercileUpperVTwoSWithinHostMaePoint{} & \numWithinHostBreakdownGroupsStabilityTercileUpperVTwoDWithinHostMaePoint{} & \numWithinHostBreakdownGroupsStabilityTercileUpperVOneSWithinHostMaePoint{} & \numWithinHostBreakdownGroupsStabilityTercileUpperVOneDWithinHostMaePoint{} \\
\addlinespace
\whgroup{(e) PBEsol(+$U$) band gap} \\
smallest-gap third & \numWithinHostBreakdownGroupsBandGapTercileLowerNHostsByMetricWithinHostMae{} & \numWithinHostBreakdownGroupsBandGapTercileLowerNSitesWithin{} & \numWithinHostBreakdownGroupsBandGapTercileLowerVTwoSZeroSkillMaePoint{} & \numWithinHostBreakdownGroupsBandGapTercileLowerVTwoSWithinHostMaePoint{} & \numWithinHostBreakdownGroupsBandGapTercileLowerVTwoDWithinHostMaePoint{} & \numWithinHostBreakdownGroupsBandGapTercileLowerVOneSWithinHostMaePoint{} & \numWithinHostBreakdownGroupsBandGapTercileLowerVOneDWithinHostMaePoint{} \\
middle third & \numWithinHostBreakdownGroupsBandGapTercileMiddleNHostsByMetricWithinHostMae{} & \numWithinHostBreakdownGroupsBandGapTercileMiddleNSitesWithin{} & \numWithinHostBreakdownGroupsBandGapTercileMiddleVTwoSZeroSkillMaePoint{} & \numWithinHostBreakdownGroupsBandGapTercileMiddleVTwoSWithinHostMaePoint{} & \numWithinHostBreakdownGroupsBandGapTercileMiddleVTwoDWithinHostMaePoint{} & \numWithinHostBreakdownGroupsBandGapTercileMiddleVOneSWithinHostMaePoint{} & \numWithinHostBreakdownGroupsBandGapTercileMiddleVOneDWithinHostMaePoint{} \\
largest-gap third & \numWithinHostBreakdownGroupsBandGapTercileUpperNHostsByMetricWithinHostMae{} & \numWithinHostBreakdownGroupsBandGapTercileUpperNSitesWithin{} & \numWithinHostBreakdownGroupsBandGapTercileUpperVTwoSZeroSkillMaePoint{} & \numWithinHostBreakdownGroupsBandGapTercileUpperVTwoSWithinHostMaePoint{} & \numWithinHostBreakdownGroupsBandGapTercileUpperVTwoDWithinHostMaePoint{} & \numWithinHostBreakdownGroupsBandGapTercileUpperVOneSWithinHostMaePoint{} & \numWithinHostBreakdownGroupsBandGapTercileUpperVOneDWithinHostMaePoint{} \\
\addlinespace
\whgroup{(f) Bader charge of the vacancy oxygen} \\
most negative O charge & \numWithinHostBreakdownGroupsBaderRankLowestNHostsByMetricWithinHostMae{} & \numWithinHostBreakdownGroupsBaderRankLowestNSitesWithin{} & \numWithinHostBreakdownGroupsBaderRankLowestVTwoSZeroSkillMaePoint{} & \numWithinHostBreakdownGroupsBaderRankLowestVTwoSWithinHostMaePoint{} & \numWithinHostBreakdownGroupsBaderRankLowestVTwoDWithinHostMaePoint{} & \numWithinHostBreakdownGroupsBaderRankLowestVOneSWithinHostMaePoint{} & \numWithinHostBreakdownGroupsBaderRankLowestVOneDWithinHostMaePoint{} \\
intermediate O charge & \numWithinHostBreakdownGroupsBaderRankMiddleNHostsByMetricWithinHostMae{} & \numWithinHostBreakdownGroupsBaderRankMiddleNSitesWithin{} & \numWithinHostBreakdownGroupsBaderRankMiddleVTwoSZeroSkillMaePoint{} & \numWithinHostBreakdownGroupsBaderRankMiddleVTwoSWithinHostMaePoint{} & \numWithinHostBreakdownGroupsBaderRankMiddleVTwoDWithinHostMaePoint{} & \numWithinHostBreakdownGroupsBaderRankMiddleVOneSWithinHostMaePoint{} & \numWithinHostBreakdownGroupsBaderRankMiddleVOneDWithinHostMaePoint{} \\
least negative O charge & \numWithinHostBreakdownGroupsBaderRankHighestNHostsByMetricWithinHostMae{} & \numWithinHostBreakdownGroupsBaderRankHighestNSitesWithin{} & \numWithinHostBreakdownGroupsBaderRankHighestVTwoSZeroSkillMaePoint{} & \numWithinHostBreakdownGroupsBaderRankHighestVTwoSWithinHostMaePoint{} & \numWithinHostBreakdownGroupsBaderRankHighestVTwoDWithinHostMaePoint{} & \numWithinHostBreakdownGroupsBaderRankHighestVOneSWithinHostMaePoint{} & \numWithinHostBreakdownGroupsBaderRankHighestVOneDWithinHostMaePoint{} \\
\bottomrule
\end{tabular}

\end{table}

\begin{figure}[htbp]
\centering
\sifig{fig_si_withinhost_groups}
\caption{Within-host residual MAE of S (blue) and D (green) at the largest budget (extension) by (\textbf{a}) cation,
(\textbf{b}) $d$-electron class, (\textbf{c}) perturbed-host-state-prone cations, (\textbf{d}) formation-energy
thirds, (\textbf{e}) band-gap thirds and (\textbf{f}) vacancy-oxygen Bader-charge rank. Grey ticks mark zero skill,
bars are two-sided host-cluster bootstrap intervals, filled circles and solid bars show the primary network (v2,
post hoc) and open circles and dashed bars the originally registered network (v1). Data:
\texttt{results/within\_host\_breakdown.json}.}
\label{fig:si_withinhost_groups}
\end{figure}
\FloatBarrier

\section{Hardware variation}
\label{si:hardware}

A trained network is a function of its data, its settings, its random seed and, on a GPU, the order in which
floating-point sums are accumulated. The aggregation over neighbouring atoms in each MEGNet block has no
deterministic GPU implementation, so the same run on two GPU models follows two slightly different training
trajectories, which can stop early at different passes and end in different weights. This section measures how
much that changes the predicted vacancy formation energies (extension; \texttt{docs/nrp/n1.md},
\texttt{results/nrp/nrp\_n1\_gpu\_repeat.json}).

The repeatability study (N1) reran the \numNrpNOneGpuRepeatCrossNodeSameClassSummaryN{} registered sweep runs of
resample \numCfgTuningAnchorResample{} at the two budgets \numCfgBudgetsHostsItemSix{} and
\numCfgBudgetsHostsItemFour{} (S and D, three seeds, originally registered network, unchanged settings) on four
GPU models, \numNrpNOneGpuRepeatNRuns{} runs in all. On one GPU model the runs are bitwise repeatable: a full
repeat on the same node reproduced every prediction
(\numNrpNOneGpuRepeatSameNodeRepeatSummaryBitwiseIdentical{} of \numNrpNOneGpuRepeatSameNodeRepeatSummaryN{} runs,
largest difference \numNrpNOneGpuRepeatSameNodeRepeatSummaryMaxAbsDiffEV{}~eV), and so did the same runs on a
second node with the same GPU model (\numNrpNOneGpuRepeatCrossNodeSameClassSummaryBitwiseIdentical{} of
\numNrpNOneGpuRepeatCrossNodeSameClassSummaryN{}). Between GPU models no run was bitwise identical, and single
predictions of one site differed by up to
\numNOneRecomputeCrossClassATenGpuFiveNrpMghpccOrgPassOneRtxAFourThousandRyGpuSixteenSdscOptiputerNetPassOneMaxAbsDiffEV{}~eV
between two classes, because a different stopping pass gives a different single network. The quantity that carries
the prediction in this work is the seed ensemble, and Table~\ref{tab:n1} gives its site MAE on each GPU model
minus that of the original sweep. At \numCfgBudgetsHostsItemSix{} training hosts the GPU model moves the ensemble
MAE of S by \numNrpNOneGpuRepeatEnsembleGeforceRtxThreeThousandNinetySSixHundredFiftyFourRyGpuTwoSdscOptiputerNetPassOneEnsembleMaeMinusRef{}
to \numNrpNOneGpuRepeatEnsembleLFourSSixHundredFiftyFourNautilusItGpuFiveFullertonEduPassOneEnsembleMaeMinusRef{}~eV
and that of D by \numNrpNOneGpuRepeatEnsembleGeforceRtxThreeThousandNinetyDStateSixHundredFiftyFourRyGpuTwoSdscOptiputerNetPassOneEnsembleMaeMinusRef{}
to \numNrpNOneGpuRepeatEnsembleRtxAFourThousandDStateSixHundredFiftyFourRyGpuSixteenSdscOptiputerNetPassOneEnsembleMaeMinusRef{}~eV;
at \numCfgBudgetsHostsItemFour{} hosts, where fewer DFT labels constrain the learned mapping, the shift of S
reaches \numNrpNOneGpuRepeatEnsembleRtxAFourThousandSTwoHundredRyGpuSixteenSdscOptiputerNetPassOneEnsembleMaeMinusRef{}~eV. The original sweep did
not record its GPU model, so its differences from the N1 runs mix GPU model with any difference in environment; the
comparisons within N1 share one code commit and one software environment.

\begin{table}[htbp]
\caption{GPU repeatability (N1; extension, descriptive). Seed-ensemble site MAE on each GPU model minus that of
the original sweep (eV), for the registered runs of resample \numCfgTuningAnchorResample{} of the originally
registered network; sweep MAE: the seed-ensemble MAE of the original sweep; seed SD: the standard deviation of the
single-run MAE over the three seeds of the original sweep, for scale. The \gpushortaten{} column holds the first pass on the
first node; the second pass and the second node are bitwise identical to it. No \gpufourty{} was free during the
study. Record: \texttt{results/nrp/nrp\_n1\_gpu\_repeat.json}, key \texttt{payload.ensemble}.}
\label{tab:n1}
\centering\small
% Generated by paper/si/tools/make_si_tables.py from results/nrp/nrp_n1_gpu_repeat.json (payload.ensemble). Do not edit.
\begin{tabular}{@{}llrrrrrr@{}}
\toprule
Model & Budget & \gpushortaten{} & \gpushortthreeoh{} & \gpushortafour{} & \gpushortlfour{} & Sweep MAE & Seed SD \\
\midrule
S & \numCfgBudgetsHostsItemSix{} & \numNrpNOneGpuRepeatEnsembleATenSSixHundredFiftyFourGpuFiveNrpMghpccOrgPassOneEnsembleMaeMinusRef{} & \numNrpNOneGpuRepeatEnsembleGeforceRtxThreeThousandNinetySSixHundredFiftyFourRyGpuTwoSdscOptiputerNetPassOneEnsembleMaeMinusRef{} & \numNrpNOneGpuRepeatEnsembleRtxAFourThousandSSixHundredFiftyFourRyGpuSixteenSdscOptiputerNetPassOneEnsembleMaeMinusRef{} & \numNrpNOneGpuRepeatEnsembleLFourSSixHundredFiftyFourNautilusItGpuFiveFullertonEduPassOneEnsembleMaeMinusRef{} & \numNrpNOneGpuRepeatEnsembleATenSSixHundredFiftyFourGpuFiveNrpMghpccOrgPassOneRefEnsembleMae{} & \numNrpNOneGpuRepeatEnsembleATenSSixHundredFiftyFourGpuFiveNrpMghpccOrgPassOneSeedSdResampleZero{} \\
S & \numCfgBudgetsHostsItemFour{} & \numNrpNOneGpuRepeatEnsembleATenSTwoHundredGpuFiveNrpMghpccOrgPassOneEnsembleMaeMinusRef{} & \numNrpNOneGpuRepeatEnsembleGeforceRtxThreeThousandNinetySTwoHundredRyGpuTwoSdscOptiputerNetPassOneEnsembleMaeMinusRef{} & \numNrpNOneGpuRepeatEnsembleRtxAFourThousandSTwoHundredRyGpuSixteenSdscOptiputerNetPassOneEnsembleMaeMinusRef{} & \numNrpNOneGpuRepeatEnsembleLFourSTwoHundredNautilusItGpuFiveFullertonEduPassOneEnsembleMaeMinusRef{} & \numNrpNOneGpuRepeatEnsembleATenSTwoHundredGpuFiveNrpMghpccOrgPassOneRefEnsembleMae{} & \numNrpNOneGpuRepeatEnsembleATenSTwoHundredGpuFiveNrpMghpccOrgPassOneSeedSdResampleZero{} \\
D & \numCfgBudgetsHostsItemSix{} & \numNrpNOneGpuRepeatEnsembleATenDStateSixHundredFiftyFourGpuFiveNrpMghpccOrgPassOneEnsembleMaeMinusRef{} & \numNrpNOneGpuRepeatEnsembleGeforceRtxThreeThousandNinetyDStateSixHundredFiftyFourRyGpuTwoSdscOptiputerNetPassOneEnsembleMaeMinusRef{} & \numNrpNOneGpuRepeatEnsembleRtxAFourThousandDStateSixHundredFiftyFourRyGpuSixteenSdscOptiputerNetPassOneEnsembleMaeMinusRef{} & \numNrpNOneGpuRepeatEnsembleLFourDStateSixHundredFiftyFourNautilusItGpuFiveFullertonEduPassOneEnsembleMaeMinusRef{} & \numNrpNOneGpuRepeatEnsembleATenDStateSixHundredFiftyFourGpuFiveNrpMghpccOrgPassOneRefEnsembleMae{} & \numNrpNOneGpuRepeatEnsembleATenDStateSixHundredFiftyFourGpuFiveNrpMghpccOrgPassOneSeedSdResampleZero{} \\
D & \numCfgBudgetsHostsItemFour{} & \numNrpNOneGpuRepeatEnsembleATenDStateTwoHundredGpuFiveNrpMghpccOrgPassOneEnsembleMaeMinusRef{} & \numNrpNOneGpuRepeatEnsembleGeforceRtxThreeThousandNinetyDStateTwoHundredRyGpuTwoSdscOptiputerNetPassOneEnsembleMaeMinusRef{} & \numNrpNOneGpuRepeatEnsembleRtxAFourThousandDStateTwoHundredRyGpuSixteenSdscOptiputerNetPassOneEnsembleMaeMinusRef{} & \numNrpNOneGpuRepeatEnsembleLFourDStateTwoHundredNautilusItGpuFiveFullertonEduPassOneEnsembleMaeMinusRef{} & \numNrpNOneGpuRepeatEnsembleATenDStateTwoHundredGpuFiveNrpMghpccOrgPassOneRefEnsembleMae{} & \numNrpNOneGpuRepeatEnsembleATenDStateTwoHundredGpuFiveNrpMghpccOrgPassOneSeedSdResampleZero{} \\
\bottomrule
\end{tabular}

\end{table}

The extensions trained on whichever GPU of the pool was free, so their runs mix GPU models. Figure~\ref{fig:gpu}
and Table~\ref{tab:gpu_class} give the advantage of the descriptors, $A = \mathrm{MAE}_\mathrm{S} -
\mathrm{MAE}_\mathrm{D}$, separately for the runs of each GPU model, using only cells in which the S run and the D
run of the same resample or fold, budget and seed ran on the same model (extension, descriptive, no test). On the
primary network (v2, post hoc) at \numCfgBudgetsHostsItemSix{} hosts the cluster classes give $A$ of
\numPrecisionABFourHundredBSixHundredFiftyFourDescriptiveGpuClassPrecisionSixHundredFiftyFourVTwoRTXFourThousandNinetyAPairedCellsA{}
to \numPrecisionABFourHundredBSixHundredFiftyFourDescriptiveGpuClassPrecisionSixHundredFiftyFourVTwoATenAPairedCellsA{}~eV,
against \numPrecisionABFourHundredBSixHundredFiftyFourDescriptiveGpuClassPrecisionSixHundredFiftyFourVTwoOriginalSweepAPairedCellsA{}~eV
for the original sweep runs; on the originally registered network the classes lie closer together. The GPU model
was not randomised with respect to resample, fold or seed, and the original sweep holds only the first resamples and
seeds, so a difference between classes cannot be separated from a difference in the resamples they ran and is not
read as an effect of the hardware. One cluster class, \gpufourty{}, is absent from N1 and ran many of the v2
precision runs. The non-inferiority statements of the article rest on the pooled runs at \numCfgBudgetsHostsItemSix{} training
hosts and at the cross-validated training sizes, on both networks; nothing in this section changes them, but every prediction carries a GPU
component of the size of Table~\ref{tab:n1} that is not separable from seed and resample variation.

\begin{figure}[htbp]
\centering
\sifig{fig_si_gpu_class}
\caption{The advantage $A$ by GPU class of the training runs (extension, descriptive, no test). $A =
\mathrm{MAE}_\mathrm{S} - \mathrm{MAE}_\mathrm{D}$ (eV) of single runs, averaged over the paired cells in which the S
run and the D run of the same resample or fold, budget and seed ran on the same GPU class. (\textbf{a}, \textbf{b})
Precision extension at the two largest budgets; (\textbf{c}, \textbf{d}) grouped $K$-fold cross-validation at the two
fold counts. The classes on the horizontal axis are the GPU models of the cluster pool, read from the run records,
and ``sweep'', the registered sweep runs of the original resamples and seeds (precision only). GPU class was not
randomised with respect to resample, fold or seed, so differences between classes are not separable from them and
are not read as hardware effects. Markers: mean over the paired cells; bars: descriptive percentile intervals of a
bootstrap over the cells; classes without a paired cell are not drawn. Filled markers and solid bars: primary network
(v2), adopted after diagnostics of the originally registered network; open markers and dashed bars: originally
registered network (v1). Diamonds: precision extension; triangles: grouped $K$-fold cross-validation. Dashed line:
the tolerance $\delta$; dotted line: zero. Values in Table~\ref{tab:gpu_class}.}
\label{fig:gpu}
\end{figure}

{\footnotesize\tighttabs
% Generated by paper/si/tools/make_si_tables.py from results/precision_a_b400_b654.json, results/cv_extended_budget.json (payload.descriptive.gpu_class). Do not edit.
\begin{longtable}{@{}lllllrrl@{}}
\caption{The advantage $A$ by GPU class of the training runs (extension, descriptive, no test). $A$ is the mean over paired cells (the S and D runs of the same resample or fold, budget and seed, both on that class) of the single-run difference MAE$_\mathrm{S}$ $-$ MAE$_\mathrm{D}$ (eV); the interval is a descriptive percentile interval of a bootstrap over the cells. Runs: all runs of S and of D on that class at that point. ``Original sweep'' holds the registered sweep runs of the original resamples and seeds. Records: \recprecision{} and \reccv{}, key \texttt{payload.descriptive.gpu\_class}.}\label{tab:gpu_class}\\
\toprule
Design & Point & Network & GPU class & Runs S / D & Cells & $A$ & Interval \\
\midrule
\endfirsthead
\toprule
Design & Point & Network & GPU class & Runs S / D & Cells & $A$ & Interval \\
\midrule
\endhead
\bottomrule
\endlastfoot
precision & \numCfgBudgetsHostsItemFive{} & v2 & \gpushortaten{} & \numPrecisionABFourHundredBSixHundredFiftyFourDescriptiveGpuClassPrecisionFourHundredVTwoATenNRunsS{} / \numPrecisionABFourHundredBSixHundredFiftyFourDescriptiveGpuClassPrecisionFourHundredVTwoATenNRunsD{} & \numPrecisionABFourHundredBSixHundredFiftyFourDescriptiveGpuClassPrecisionFourHundredVTwoATenNPairedCells{} & \numPrecisionABFourHundredBSixHundredFiftyFourDescriptiveGpuClassPrecisionFourHundredVTwoATenAPairedCellsA{} & \ci{\numPrecisionABFourHundredBSixHundredFiftyFourDescriptiveGpuClassPrecisionFourHundredVTwoATenAPairedCellsCiNinetyFiveDescriptiveItemOne{}}{\numPrecisionABFourHundredBSixHundredFiftyFourDescriptiveGpuClassPrecisionFourHundredVTwoATenAPairedCellsCiNinetyFiveDescriptiveItemTwo{}} \\
 &  & v2 & \gpushortthreeoh{} & \numPrecisionABFourHundredBSixHundredFiftyFourDescriptiveGpuClassPrecisionFourHundredVTwoRTXThreeThousandNinetyNRunsS{} / \numPrecisionABFourHundredBSixHundredFiftyFourDescriptiveGpuClassPrecisionFourHundredVTwoRTXThreeThousandNinetyNRunsD{} & \numPrecisionABFourHundredBSixHundredFiftyFourDescriptiveGpuClassPrecisionFourHundredVTwoRTXThreeThousandNinetyNPairedCells{} & \numPrecisionABFourHundredBSixHundredFiftyFourDescriptiveGpuClassPrecisionFourHundredVTwoRTXThreeThousandNinetyAPairedCellsA{} & \ci{\numPrecisionABFourHundredBSixHundredFiftyFourDescriptiveGpuClassPrecisionFourHundredVTwoRTXThreeThousandNinetyAPairedCellsCiNinetyFiveDescriptiveItemOne{}}{\numPrecisionABFourHundredBSixHundredFiftyFourDescriptiveGpuClassPrecisionFourHundredVTwoRTXThreeThousandNinetyAPairedCellsCiNinetyFiveDescriptiveItemTwo{}} \\
 &  & v2 & \gpushortlfour{} & \numPrecisionABFourHundredBSixHundredFiftyFourDescriptiveGpuClassPrecisionFourHundredVTwoLFourNRunsS{} / \numPrecisionABFourHundredBSixHundredFiftyFourDescriptiveGpuClassPrecisionFourHundredVTwoLFourNRunsD{} & \numPrecisionABFourHundredBSixHundredFiftyFourDescriptiveGpuClassPrecisionFourHundredVTwoLFourNPairedCells{} & \numPrecisionABFourHundredBSixHundredFiftyFourDescriptiveGpuClassPrecisionFourHundredVTwoLFourAPairedCellsA{} & \ci{\numPrecisionABFourHundredBSixHundredFiftyFourDescriptiveGpuClassPrecisionFourHundredVTwoLFourAPairedCellsCiNinetyFiveDescriptiveItemOne{}}{\numPrecisionABFourHundredBSixHundredFiftyFourDescriptiveGpuClassPrecisionFourHundredVTwoLFourAPairedCellsCiNinetyFiveDescriptiveItemTwo{}} \\
 &  & v2 & \gpushortfourty{} (not in N1) & \numPrecisionABFourHundredBSixHundredFiftyFourDescriptiveGpuClassPrecisionFourHundredVTwoRTXFourThousandNinetyNRunsS{} / \numPrecisionABFourHundredBSixHundredFiftyFourDescriptiveGpuClassPrecisionFourHundredVTwoRTXFourThousandNinetyNRunsD{} & \numPrecisionABFourHundredBSixHundredFiftyFourDescriptiveGpuClassPrecisionFourHundredVTwoRTXFourThousandNinetyNPairedCells{} & \numPrecisionABFourHundredBSixHundredFiftyFourDescriptiveGpuClassPrecisionFourHundredVTwoRTXFourThousandNinetyAPairedCellsA{} & \ci{\numPrecisionABFourHundredBSixHundredFiftyFourDescriptiveGpuClassPrecisionFourHundredVTwoRTXFourThousandNinetyAPairedCellsCiNinetyFiveDescriptiveItemOne{}}{\numPrecisionABFourHundredBSixHundredFiftyFourDescriptiveGpuClassPrecisionFourHundredVTwoRTXFourThousandNinetyAPairedCellsCiNinetyFiveDescriptiveItemTwo{}} \\
 &  & v2 & original sweep & \numPrecisionABFourHundredBSixHundredFiftyFourDescriptiveGpuClassPrecisionFourHundredVTwoOriginalSweepNRunsS{} / \numPrecisionABFourHundredBSixHundredFiftyFourDescriptiveGpuClassPrecisionFourHundredVTwoOriginalSweepNRunsD{} & \numPrecisionABFourHundredBSixHundredFiftyFourDescriptiveGpuClassPrecisionFourHundredVTwoOriginalSweepNPairedCells{} & \numPrecisionABFourHundredBSixHundredFiftyFourDescriptiveGpuClassPrecisionFourHundredVTwoOriginalSweepAPairedCellsA{} & \ci{\numPrecisionABFourHundredBSixHundredFiftyFourDescriptiveGpuClassPrecisionFourHundredVTwoOriginalSweepAPairedCellsCiNinetyFiveDescriptiveItemOne{}}{\numPrecisionABFourHundredBSixHundredFiftyFourDescriptiveGpuClassPrecisionFourHundredVTwoOriginalSweepAPairedCellsCiNinetyFiveDescriptiveItemTwo{}} \\
 &  & v1 & \gpushortaten{} & \numPrecisionABFourHundredBSixHundredFiftyFourDescriptiveGpuClassPrecisionFourHundredVOneATenNRunsS{} / \numPrecisionABFourHundredBSixHundredFiftyFourDescriptiveGpuClassPrecisionFourHundredVOneATenNRunsD{} & \numPrecisionABFourHundredBSixHundredFiftyFourDescriptiveGpuClassPrecisionFourHundredVOneATenNPairedCells{} & \numPrecisionABFourHundredBSixHundredFiftyFourDescriptiveGpuClassPrecisionFourHundredVOneATenAPairedCellsA{} & \ci{\numPrecisionABFourHundredBSixHundredFiftyFourDescriptiveGpuClassPrecisionFourHundredVOneATenAPairedCellsCiNinetyFiveDescriptiveItemOne{}}{\numPrecisionABFourHundredBSixHundredFiftyFourDescriptiveGpuClassPrecisionFourHundredVOneATenAPairedCellsCiNinetyFiveDescriptiveItemTwo{}} \\
 &  & v1 & \gpushortthreeoh{} & \numPrecisionABFourHundredBSixHundredFiftyFourDescriptiveGpuClassPrecisionFourHundredVOneRTXThreeThousandNinetyNRunsS{} / \numPrecisionABFourHundredBSixHundredFiftyFourDescriptiveGpuClassPrecisionFourHundredVOneRTXThreeThousandNinetyNRunsD{} & \numPrecisionABFourHundredBSixHundredFiftyFourDescriptiveGpuClassPrecisionFourHundredVOneRTXThreeThousandNinetyNPairedCells{} & \numPrecisionABFourHundredBSixHundredFiftyFourDescriptiveGpuClassPrecisionFourHundredVOneRTXThreeThousandNinetyAPairedCellsA{} & \ci{\numPrecisionABFourHundredBSixHundredFiftyFourDescriptiveGpuClassPrecisionFourHundredVOneRTXThreeThousandNinetyAPairedCellsCiNinetyFiveDescriptiveItemOne{}}{\numPrecisionABFourHundredBSixHundredFiftyFourDescriptiveGpuClassPrecisionFourHundredVOneRTXThreeThousandNinetyAPairedCellsCiNinetyFiveDescriptiveItemTwo{}} \\
 &  & v1 & \gpushortlfour{} & \numPrecisionABFourHundredBSixHundredFiftyFourDescriptiveGpuClassPrecisionFourHundredVOneLFourNRunsS{} / \numPrecisionABFourHundredBSixHundredFiftyFourDescriptiveGpuClassPrecisionFourHundredVOneLFourNRunsD{} & \numPrecisionABFourHundredBSixHundredFiftyFourDescriptiveGpuClassPrecisionFourHundredVOneLFourNPairedCells{} & \numPrecisionABFourHundredBSixHundredFiftyFourDescriptiveGpuClassPrecisionFourHundredVOneLFourAPairedCellsA{} & \ci{\numPrecisionABFourHundredBSixHundredFiftyFourDescriptiveGpuClassPrecisionFourHundredVOneLFourAPairedCellsCiNinetyFiveDescriptiveItemOne{}}{\numPrecisionABFourHundredBSixHundredFiftyFourDescriptiveGpuClassPrecisionFourHundredVOneLFourAPairedCellsCiNinetyFiveDescriptiveItemTwo{}} \\
 &  & v1 & \gpushortfourty{} (not in N1) & \numPrecisionABFourHundredBSixHundredFiftyFourDescriptiveGpuClassPrecisionFourHundredVOneRTXFourThousandNinetyNRunsS{} / \numPrecisionABFourHundredBSixHundredFiftyFourDescriptiveGpuClassPrecisionFourHundredVOneRTXFourThousandNinetyNRunsD{} & \numPrecisionABFourHundredBSixHundredFiftyFourDescriptiveGpuClassPrecisionFourHundredVOneRTXFourThousandNinetyNPairedCells{} & -- &  \\
 &  & v1 & original sweep & \numPrecisionABFourHundredBSixHundredFiftyFourDescriptiveGpuClassPrecisionFourHundredVOneOriginalSweepNRunsS{} / \numPrecisionABFourHundredBSixHundredFiftyFourDescriptiveGpuClassPrecisionFourHundredVOneOriginalSweepNRunsD{} & \numPrecisionABFourHundredBSixHundredFiftyFourDescriptiveGpuClassPrecisionFourHundredVOneOriginalSweepNPairedCells{} & \numPrecisionABFourHundredBSixHundredFiftyFourDescriptiveGpuClassPrecisionFourHundredVOneOriginalSweepAPairedCellsA{} & \ci{\numPrecisionABFourHundredBSixHundredFiftyFourDescriptiveGpuClassPrecisionFourHundredVOneOriginalSweepAPairedCellsCiNinetyFiveDescriptiveItemOne{}}{\numPrecisionABFourHundredBSixHundredFiftyFourDescriptiveGpuClassPrecisionFourHundredVOneOriginalSweepAPairedCellsCiNinetyFiveDescriptiveItemTwo{}} \\
\addlinespace
precision & \numCfgBudgetsHostsItemSix{} & v2 & \gpushortaten{} & \numPrecisionABFourHundredBSixHundredFiftyFourDescriptiveGpuClassPrecisionSixHundredFiftyFourVTwoATenNRunsS{} / \numPrecisionABFourHundredBSixHundredFiftyFourDescriptiveGpuClassPrecisionSixHundredFiftyFourVTwoATenNRunsD{} & \numPrecisionABFourHundredBSixHundredFiftyFourDescriptiveGpuClassPrecisionSixHundredFiftyFourVTwoATenNPairedCells{} & \numPrecisionABFourHundredBSixHundredFiftyFourDescriptiveGpuClassPrecisionSixHundredFiftyFourVTwoATenAPairedCellsA{} & \ci{\numPrecisionABFourHundredBSixHundredFiftyFourDescriptiveGpuClassPrecisionSixHundredFiftyFourVTwoATenAPairedCellsCiNinetyFiveDescriptiveItemOne{}}{\numPrecisionABFourHundredBSixHundredFiftyFourDescriptiveGpuClassPrecisionSixHundredFiftyFourVTwoATenAPairedCellsCiNinetyFiveDescriptiveItemTwo{}} \\
 &  & v2 & \gpushortthreeoh{} & \numPrecisionABFourHundredBSixHundredFiftyFourDescriptiveGpuClassPrecisionSixHundredFiftyFourVTwoRTXThreeThousandNinetyNRunsS{} / \numPrecisionABFourHundredBSixHundredFiftyFourDescriptiveGpuClassPrecisionSixHundredFiftyFourVTwoRTXThreeThousandNinetyNRunsD{} & \numPrecisionABFourHundredBSixHundredFiftyFourDescriptiveGpuClassPrecisionSixHundredFiftyFourVTwoRTXThreeThousandNinetyNPairedCells{} & \numPrecisionABFourHundredBSixHundredFiftyFourDescriptiveGpuClassPrecisionSixHundredFiftyFourVTwoRTXThreeThousandNinetyAPairedCellsA{} & \ci{\numPrecisionABFourHundredBSixHundredFiftyFourDescriptiveGpuClassPrecisionSixHundredFiftyFourVTwoRTXThreeThousandNinetyAPairedCellsCiNinetyFiveDescriptiveItemOne{}}{\numPrecisionABFourHundredBSixHundredFiftyFourDescriptiveGpuClassPrecisionSixHundredFiftyFourVTwoRTXThreeThousandNinetyAPairedCellsCiNinetyFiveDescriptiveItemTwo{}} \\
 &  & v2 & \gpushortlfour{} & \numPrecisionABFourHundredBSixHundredFiftyFourDescriptiveGpuClassPrecisionSixHundredFiftyFourVTwoLFourNRunsS{} / \numPrecisionABFourHundredBSixHundredFiftyFourDescriptiveGpuClassPrecisionSixHundredFiftyFourVTwoLFourNRunsD{} & \numPrecisionABFourHundredBSixHundredFiftyFourDescriptiveGpuClassPrecisionSixHundredFiftyFourVTwoLFourNPairedCells{} & -- &  \\
 &  & v2 & \gpushortfourty{} (not in N1) & \numPrecisionABFourHundredBSixHundredFiftyFourDescriptiveGpuClassPrecisionSixHundredFiftyFourVTwoRTXFourThousandNinetyNRunsS{} / \numPrecisionABFourHundredBSixHundredFiftyFourDescriptiveGpuClassPrecisionSixHundredFiftyFourVTwoRTXFourThousandNinetyNRunsD{} & \numPrecisionABFourHundredBSixHundredFiftyFourDescriptiveGpuClassPrecisionSixHundredFiftyFourVTwoRTXFourThousandNinetyNPairedCells{} & \numPrecisionABFourHundredBSixHundredFiftyFourDescriptiveGpuClassPrecisionSixHundredFiftyFourVTwoRTXFourThousandNinetyAPairedCellsA{} & \ci{\numPrecisionABFourHundredBSixHundredFiftyFourDescriptiveGpuClassPrecisionSixHundredFiftyFourVTwoRTXFourThousandNinetyAPairedCellsCiNinetyFiveDescriptiveItemOne{}}{\numPrecisionABFourHundredBSixHundredFiftyFourDescriptiveGpuClassPrecisionSixHundredFiftyFourVTwoRTXFourThousandNinetyAPairedCellsCiNinetyFiveDescriptiveItemTwo{}} \\
 &  & v2 & original sweep & \numPrecisionABFourHundredBSixHundredFiftyFourDescriptiveGpuClassPrecisionSixHundredFiftyFourVTwoOriginalSweepNRunsS{} / \numPrecisionABFourHundredBSixHundredFiftyFourDescriptiveGpuClassPrecisionSixHundredFiftyFourVTwoOriginalSweepNRunsD{} & \numPrecisionABFourHundredBSixHundredFiftyFourDescriptiveGpuClassPrecisionSixHundredFiftyFourVTwoOriginalSweepNPairedCells{} & \numPrecisionABFourHundredBSixHundredFiftyFourDescriptiveGpuClassPrecisionSixHundredFiftyFourVTwoOriginalSweepAPairedCellsA{} & \ci{\numPrecisionABFourHundredBSixHundredFiftyFourDescriptiveGpuClassPrecisionSixHundredFiftyFourVTwoOriginalSweepAPairedCellsCiNinetyFiveDescriptiveItemOne{}}{\numPrecisionABFourHundredBSixHundredFiftyFourDescriptiveGpuClassPrecisionSixHundredFiftyFourVTwoOriginalSweepAPairedCellsCiNinetyFiveDescriptiveItemTwo{}} \\
 &  & v1 & \gpushortaten{} & \numPrecisionABFourHundredBSixHundredFiftyFourDescriptiveGpuClassPrecisionSixHundredFiftyFourVOneATenNRunsS{} / \numPrecisionABFourHundredBSixHundredFiftyFourDescriptiveGpuClassPrecisionSixHundredFiftyFourVOneATenNRunsD{} & \numPrecisionABFourHundredBSixHundredFiftyFourDescriptiveGpuClassPrecisionSixHundredFiftyFourVOneATenNPairedCells{} & \numPrecisionABFourHundredBSixHundredFiftyFourDescriptiveGpuClassPrecisionSixHundredFiftyFourVOneATenAPairedCellsA{} & \ci{\numPrecisionABFourHundredBSixHundredFiftyFourDescriptiveGpuClassPrecisionSixHundredFiftyFourVOneATenAPairedCellsCiNinetyFiveDescriptiveItemOne{}}{\numPrecisionABFourHundredBSixHundredFiftyFourDescriptiveGpuClassPrecisionSixHundredFiftyFourVOneATenAPairedCellsCiNinetyFiveDescriptiveItemTwo{}} \\
 &  & v1 & \gpushortthreeoh{} & \numPrecisionABFourHundredBSixHundredFiftyFourDescriptiveGpuClassPrecisionSixHundredFiftyFourVOneRTXThreeThousandNinetyNRunsS{} / \numPrecisionABFourHundredBSixHundredFiftyFourDescriptiveGpuClassPrecisionSixHundredFiftyFourVOneRTXThreeThousandNinetyNRunsD{} & \numPrecisionABFourHundredBSixHundredFiftyFourDescriptiveGpuClassPrecisionSixHundredFiftyFourVOneRTXThreeThousandNinetyNPairedCells{} & \numPrecisionABFourHundredBSixHundredFiftyFourDescriptiveGpuClassPrecisionSixHundredFiftyFourVOneRTXThreeThousandNinetyAPairedCellsA{} & \ci{\numPrecisionABFourHundredBSixHundredFiftyFourDescriptiveGpuClassPrecisionSixHundredFiftyFourVOneRTXThreeThousandNinetyAPairedCellsCiNinetyFiveDescriptiveItemOne{}}{\numPrecisionABFourHundredBSixHundredFiftyFourDescriptiveGpuClassPrecisionSixHundredFiftyFourVOneRTXThreeThousandNinetyAPairedCellsCiNinetyFiveDescriptiveItemTwo{}} \\
 &  & v1 & \gpushortlfour{} & \numPrecisionABFourHundredBSixHundredFiftyFourDescriptiveGpuClassPrecisionSixHundredFiftyFourVOneLFourNRunsS{} / \numPrecisionABFourHundredBSixHundredFiftyFourDescriptiveGpuClassPrecisionSixHundredFiftyFourVOneLFourNRunsD{} & \numPrecisionABFourHundredBSixHundredFiftyFourDescriptiveGpuClassPrecisionSixHundredFiftyFourVOneLFourNPairedCells{} & \numPrecisionABFourHundredBSixHundredFiftyFourDescriptiveGpuClassPrecisionSixHundredFiftyFourVOneLFourAPairedCellsA{} & \ci{\numPrecisionABFourHundredBSixHundredFiftyFourDescriptiveGpuClassPrecisionSixHundredFiftyFourVOneLFourAPairedCellsCiNinetyFiveDescriptiveItemOne{}}{\numPrecisionABFourHundredBSixHundredFiftyFourDescriptiveGpuClassPrecisionSixHundredFiftyFourVOneLFourAPairedCellsCiNinetyFiveDescriptiveItemTwo{}} \\
 &  & v1 & \gpushortfourty{} (not in N1) & \numPrecisionABFourHundredBSixHundredFiftyFourDescriptiveGpuClassPrecisionSixHundredFiftyFourVOneRTXFourThousandNinetyNRunsS{} / \numPrecisionABFourHundredBSixHundredFiftyFourDescriptiveGpuClassPrecisionSixHundredFiftyFourVOneRTXFourThousandNinetyNRunsD{} & \numPrecisionABFourHundredBSixHundredFiftyFourDescriptiveGpuClassPrecisionSixHundredFiftyFourVOneRTXFourThousandNinetyNPairedCells{} & \numPrecisionABFourHundredBSixHundredFiftyFourDescriptiveGpuClassPrecisionSixHundredFiftyFourVOneRTXFourThousandNinetyAPairedCellsA{} & \ci{\numPrecisionABFourHundredBSixHundredFiftyFourDescriptiveGpuClassPrecisionSixHundredFiftyFourVOneRTXFourThousandNinetyAPairedCellsCiNinetyFiveDescriptiveItemOne{}}{\numPrecisionABFourHundredBSixHundredFiftyFourDescriptiveGpuClassPrecisionSixHundredFiftyFourVOneRTXFourThousandNinetyAPairedCellsCiNinetyFiveDescriptiveItemTwo{}} \\
 &  & v1 & original sweep & \numPrecisionABFourHundredBSixHundredFiftyFourDescriptiveGpuClassPrecisionSixHundredFiftyFourVOneOriginalSweepNRunsS{} / \numPrecisionABFourHundredBSixHundredFiftyFourDescriptiveGpuClassPrecisionSixHundredFiftyFourVOneOriginalSweepNRunsD{} & \numPrecisionABFourHundredBSixHundredFiftyFourDescriptiveGpuClassPrecisionSixHundredFiftyFourVOneOriginalSweepNPairedCells{} & \numPrecisionABFourHundredBSixHundredFiftyFourDescriptiveGpuClassPrecisionSixHundredFiftyFourVOneOriginalSweepAPairedCellsA{} & \ci{\numPrecisionABFourHundredBSixHundredFiftyFourDescriptiveGpuClassPrecisionSixHundredFiftyFourVOneOriginalSweepAPairedCellsCiNinetyFiveDescriptiveItemOne{}}{\numPrecisionABFourHundredBSixHundredFiftyFourDescriptiveGpuClassPrecisionSixHundredFiftyFourVOneOriginalSweepAPairedCellsCiNinetyFiveDescriptiveItemTwo{}} \\
\addlinespace
grouped CV & $K = \numCvExtendedBudgetDescriptiveGpuClassCvTenVTwoATenPoint{}$ & v2 & \gpushortaten{} & \numCvExtendedBudgetDescriptiveGpuClassCvTenVTwoATenNRunsS{} / \numCvExtendedBudgetDescriptiveGpuClassCvTenVTwoATenNRunsD{} & \numCvExtendedBudgetDescriptiveGpuClassCvTenVTwoATenNPairedCells{} & \numCvExtendedBudgetDescriptiveGpuClassCvTenVTwoATenAPairedCellsA{} & \ci{\numCvExtendedBudgetDescriptiveGpuClassCvTenVTwoATenAPairedCellsCiNinetyFiveDescriptiveItemOne{}}{\numCvExtendedBudgetDescriptiveGpuClassCvTenVTwoATenAPairedCellsCiNinetyFiveDescriptiveItemTwo{}} \\
 &  & v2 & \gpushortthreeoh{} & \numCvExtendedBudgetDescriptiveGpuClassCvTenVTwoRTXThreeThousandNinetyNRunsS{} / \numCvExtendedBudgetDescriptiveGpuClassCvTenVTwoRTXThreeThousandNinetyNRunsD{} & \numCvExtendedBudgetDescriptiveGpuClassCvTenVTwoRTXThreeThousandNinetyNPairedCells{} & \numCvExtendedBudgetDescriptiveGpuClassCvTenVTwoRTXThreeThousandNinetyAPairedCellsA{} & \ci{\numCvExtendedBudgetDescriptiveGpuClassCvTenVTwoRTXThreeThousandNinetyAPairedCellsCiNinetyFiveDescriptiveItemOne{}}{\numCvExtendedBudgetDescriptiveGpuClassCvTenVTwoRTXThreeThousandNinetyAPairedCellsCiNinetyFiveDescriptiveItemTwo{}} \\
 &  & v1 & \gpushortaten{} & \numCvExtendedBudgetDescriptiveGpuClassCvTenVOneATenNRunsS{} / \numCvExtendedBudgetDescriptiveGpuClassCvTenVOneATenNRunsD{} & \numCvExtendedBudgetDescriptiveGpuClassCvTenVOneATenNPairedCells{} & \numCvExtendedBudgetDescriptiveGpuClassCvTenVOneATenAPairedCellsA{} & \ci{\numCvExtendedBudgetDescriptiveGpuClassCvTenVOneATenAPairedCellsCiNinetyFiveDescriptiveItemOne{}}{\numCvExtendedBudgetDescriptiveGpuClassCvTenVOneATenAPairedCellsCiNinetyFiveDescriptiveItemTwo{}} \\
 &  & v1 & \gpushortlfour{} & \numCvExtendedBudgetDescriptiveGpuClassCvTenVOneLFourNRunsS{} / \numCvExtendedBudgetDescriptiveGpuClassCvTenVOneLFourNRunsD{} & \numCvExtendedBudgetDescriptiveGpuClassCvTenVOneLFourNPairedCells{} & \numCvExtendedBudgetDescriptiveGpuClassCvTenVOneLFourAPairedCellsA{} & \ci{\numCvExtendedBudgetDescriptiveGpuClassCvTenVOneLFourAPairedCellsCiNinetyFiveDescriptiveItemOne{}}{\numCvExtendedBudgetDescriptiveGpuClassCvTenVOneLFourAPairedCellsCiNinetyFiveDescriptiveItemTwo{}} \\
\addlinespace
grouped CV & $K = \numCvExtendedBudgetDescriptiveGpuClassCvTwentyVTwoATenPoint{}$ & v2 & \gpushortaten{} & \numCvExtendedBudgetDescriptiveGpuClassCvTwentyVTwoATenNRunsS{} / \numCvExtendedBudgetDescriptiveGpuClassCvTwentyVTwoATenNRunsD{} & \numCvExtendedBudgetDescriptiveGpuClassCvTwentyVTwoATenNPairedCells{} & \numCvExtendedBudgetDescriptiveGpuClassCvTwentyVTwoATenAPairedCellsA{} & \ci{\numCvExtendedBudgetDescriptiveGpuClassCvTwentyVTwoATenAPairedCellsCiNinetyFiveDescriptiveItemOne{}}{\numCvExtendedBudgetDescriptiveGpuClassCvTwentyVTwoATenAPairedCellsCiNinetyFiveDescriptiveItemTwo{}} \\
 &  & v2 & \gpushortthreeoh{} & \numCvExtendedBudgetDescriptiveGpuClassCvTwentyVTwoRTXThreeThousandNinetyNRunsS{} / \numCvExtendedBudgetDescriptiveGpuClassCvTwentyVTwoRTXThreeThousandNinetyNRunsD{} & \numCvExtendedBudgetDescriptiveGpuClassCvTwentyVTwoRTXThreeThousandNinetyNPairedCells{} & \numCvExtendedBudgetDescriptiveGpuClassCvTwentyVTwoRTXThreeThousandNinetyAPairedCellsA{} & \ci{\numCvExtendedBudgetDescriptiveGpuClassCvTwentyVTwoRTXThreeThousandNinetyAPairedCellsCiNinetyFiveDescriptiveItemOne{}}{\numCvExtendedBudgetDescriptiveGpuClassCvTwentyVTwoRTXThreeThousandNinetyAPairedCellsCiNinetyFiveDescriptiveItemTwo{}} \\
 &  & v1 & \gpushortaten{} & \numCvExtendedBudgetDescriptiveGpuClassCvTwentyVOneATenNRunsS{} / \numCvExtendedBudgetDescriptiveGpuClassCvTwentyVOneATenNRunsD{} & \numCvExtendedBudgetDescriptiveGpuClassCvTwentyVOneATenNPairedCells{} & \numCvExtendedBudgetDescriptiveGpuClassCvTwentyVOneATenAPairedCellsA{} & \ci{\numCvExtendedBudgetDescriptiveGpuClassCvTwentyVOneATenAPairedCellsCiNinetyFiveDescriptiveItemOne{}}{\numCvExtendedBudgetDescriptiveGpuClassCvTwentyVOneATenAPairedCellsCiNinetyFiveDescriptiveItemTwo{}} \\
 &  & v1 & \gpushortlfour{} & \numCvExtendedBudgetDescriptiveGpuClassCvTwentyVOneLFourNRunsS{} / \numCvExtendedBudgetDescriptiveGpuClassCvTwentyVOneLFourNRunsD{} & \numCvExtendedBudgetDescriptiveGpuClassCvTwentyVOneLFourNPairedCells{} & \numCvExtendedBudgetDescriptiveGpuClassCvTwentyVOneLFourAPairedCellsA{} & \ci{\numCvExtendedBudgetDescriptiveGpuClassCvTwentyVOneLFourAPairedCellsCiNinetyFiveDescriptiveItemOne{}}{\numCvExtendedBudgetDescriptiveGpuClassCvTwentyVOneLFourAPairedCellsCiNinetyFiveDescriptiveItemTwo{}} \\
\end{longtable}
}
\FloatBarrier

\section{Sensitivity analyses}
\label{si:sensitivities}

Three choices fixed by the registered design could, in principle, move the bound on the advantage of the electronic
descriptors, $A = \mathrm{MAE}_\mathrm{S} - \mathrm{MAE}_\mathrm{D}$, at the largest training set: the inclusion of
vacancies that leave a magnetic moment on the cell, the cap on training length and the choice of architecture
(Figure~\ref{fig:sens}; \texttt{docs/a5\_notes.md}). All three were run for the originally registered network
(v1) only; no v2 sensitivity runs exist, so the v2 statement rests on its primary result alone (an upper bound of
\numSensCappedPostHocVTwoPrimarySixHundredFiftyFourPrimaryAUpperOneSided{}~eV at \numCfgBudgetsHostsItemSix{}
hosts, post hoc).

The database treats every host as non-magnetic, but a few neutral vacancies leave a moment of more than
\numSensHighMomentVOneThresholdMuB{}~$\mu_\mathrm{B}$ on the cell, as two electrons localised on reducible cations
would; such polaron-like configurations are the ones that a non-magnetic target describes least well. Retraining S
and D with these entries removed from training and test (extension) changes $A$ on the common test sites from
\numSensHighMomentVOnePrimaryCommonSitesAMean{} to \numSensHighMomentVOneHighMomentAMean{}~eV, a paired change of
\numSensHighMomentVOneChangeHighMomentMinusPrimaryAMean{}~eV
\ci{\numSensHighMomentVOneChangeHighMomentMinusPrimaryACiLo{}}{\numSensHighMomentVOneChangeHighMomentMinusPrimaryACiHi{}},
and the one-sided upper bound moves from \numSensHighMomentVOnePrimaryCommonSitesAUpperOneSided{} to
\numSensHighMomentVOneHighMomentAUpperOneSided{}~eV, below the tolerance in both cases. The excluded entries are
rare in any test set, so the exclusion cannot move the pooled bound by much, and it does not.

Every primary run that stopped at the cap of \numCfgTrainingMaxEpochs{} passes was rerun with a cap of
\numCfgSensitivityEpochCapMaxEpochs{} passes, and the capped run was replaced by its twin (post hoc, rule logged
before the reruns). At \numCfgBudgetsHostsItemSix{} hosts $A$ moves from
\numSensCappedPostHocPerBudgetSixHundredFiftyFourPrimaryAMean{} to
\numSensCappedPostHocPerBudgetSixHundredFiftyFourCappedTwinsAMean{}~eV (paired change
\numSensCappedPostHocPerBudgetSixHundredFiftyFourChangeTwinMinusPrimaryAMean{}~eV
\ci{\numSensCappedPostHocPerBudgetSixHundredFiftyFourChangeTwinMinusPrimaryACiLo{}}{\numSensCappedPostHocPerBudgetSixHundredFiftyFourChangeTwinMinusPrimaryACiHi{}}),
and its upper bound from \numSensCappedPostHocPerBudgetSixHundredFiftyFourPrimaryAUpperOneSided{} to
\numSensCappedPostHocPerBudgetSixHundredFiftyFourCappedTwinsAUpperOneSided{}~eV, still below the tolerance. At
\numCfgBudgetsHostsItemFive{} hosts the bound stays above it under both versions
(\numSensCappedPostHocPerBudgetFourHundredPrimaryAUpperOneSided{} and
\numSensCappedPostHocPerBudgetFourHundredCappedTwinsAUpperOneSided{}~eV), so the smallest budget at which the
descriptors are bounded below the tolerance remains \numCfgBudgetsHostsItemSix{}. The twins cover only the runs of
the registered sweep; the capped runs of the extensions (Table~\ref{tab:capped}) were not rerun.

The architecture check trained an untuned CGCNN \citep{Xie2018crystal}, with fixed optimiser settings, on one outer
resample at three budgets (extension, descriptive). Its advantage falls with the number of training hosts as that of
both MEGNet networks does, from \numSensCgcnnTableCgcnnFiftyAPoint{}~eV at \numCfgBudgetsHostsItemTwo{} hosts to
\numSensCgcnnTableCgcnnSixHundredFiftyFourAPoint{}~eV
\ci{\numSensCgcnnTableCgcnnSixHundredFiftyFourACiNinetyFiveItemOne{}}{\numSensCgcnnTableCgcnnSixHundredFiftyFourACiNinetyFiveItemTwo{}}
at \numCfgBudgetsHostsItemSix{}, against \numSensCgcnnTableMegnetVOneSixHundredFiftyFourAPoint{}~eV for MEGNet v1
and \numSensCgcnnTableMegnetVTwoSixHundredFiftyFourAPoint{}~eV for MEGNet v2 on the same draws. CGCNN has
\numSensCgcnnParameterCountsCgcnnCgcnnSSixHundredFiftyFour{} trainable parameters in S, fewer than any of the tuned MEGNets
(\numSensCgcnnParameterCountsMegnetVTwoSTwoHundred{} to \numSensCgcnnParameterCountsMegnetVOneSTwoHundred{}). With one resample, the intervals reflect only the sampling of its test hosts, not split or retraining
variation, so the check shows that the fall of $A$ with training data is not peculiar to MEGNet; it does not
establish a bound for other architectures.

\begin{figure}[htbp]
\centering
\sifig{fig_si_sensitivities}
\caption{The advantage $A$ under the high-moment exclusion and the longer training cap, and an untuned CGCNN beside
the MEGNet networks. (\textbf{a}) $A = \mathrm{MAE}_\mathrm{S} - \mathrm{MAE}_\mathrm{D}$ (eV; positive when the DFT
descriptors lower the site error) at the largest training budget, from the seed-ensemble site MAE averaged over the
registered outer resamples. ``Primary'': the primary result of each network. ``High moment'' (extension): S and D
retrained with the entries whose defect magnetic moment exceeds the registered threshold removed from training and
test sites, compared on the common test sites. ``Capped twins'' (post hoc): every primary run that stopped at the
training cap replaced by its rerun with the longer cap. Both sensitivities exist for the originally registered
network only. Markers: point estimates; bars: two-sided intervals of the hierarchical bootstrap of the primary
analysis; short horizontal bars: one-sided upper bounds; dashed line: the tolerance $\delta$; dotted line: zero.
(\textbf{b}) Seed-ensemble test MAE (eV) of S (blue) and D (green) on a single outer resample at three budgets, for the
MEGNet networks (lines with bands) and for CGCNN with fixed, untuned optimiser settings (crosses, drawn just right of
their budget). Bands and bars in (b) are two-sided intervals of a host-cluster bootstrap over the test hosts of that
resample; with one resample they reflect test-host sampling only, and no test is made. Filled markers and solid
lines: primary network (v2), adopted after diagnostics of the originally registered network; open markers and
dashed lines: originally registered network (v1). Records: \texttt{results/sens\_high\_moment\_v1.json},
\texttt{results/sens\_capped\_post\_hoc.json}, \texttt{results/sens\_cgcnn.json}.}
\label{fig:sens}
\end{figure}
\FloatBarrier

\section{Latent probe: does training teach the structure-only network the electronic structure?}
\label{si:probe}

If learning formation energies from structure alone made the network learn the electronic structure of the host,
a linear map from its learned description of the vacancy site and its host to the DFT descriptors should improve as
the network sees more training oxides, and should do better than the same map applied to an untrained network
reading the same geometry. This is claim C4, registered on the originally registered network (v1;
\texttt{docs/a4\_notes.md}). For every resample and budget, the readout input of the seed-0 S checkpoint (the
vacancy-site vector, the pooled host vector and the global state) was mapped by a ridge regression, fitted on the
training sites with its regularisation chosen by cross-validation grouped by host, to each descriptor given to D,
and scored by the coefficient of determination \Rsq{} on the test sites. The class means of the host-electronic and
the site-electronic descriptors are the quantities of the rule.

The registered rule (\texttt{docs/deviations.md}, the latent-probe row of step A4, committed before any v1 probe
number existed) has two conditions for each class. The trend condition asks that the recovered fraction of
descriptor variance rise with the number of training hosts: the Spearman correlation between the logarithm of the
budget and the class-mean \Rsq{} over all \numProbeCFourVOneRuleHostTrendNCells{} resample-budget cells must have an
interval above zero. The control condition asks that, at every budget from \numCfgProbePassRuleBeatsControlsMinBudget{}
hosts on, the trained probe exceed both the same probe fitted to descriptor rows permuted among the training
hosts (shuffled labels) and the probe on an untrained network of the same architecture (the registered
random-initialisation control), each difference with an interval above zero.

The trend condition holds for both classes, with Spearman correlations of \numProbeCFourVOneRuleHostTrendRho{}
\ci{\numProbeCFourVOneRuleHostTrendCiNinetyFiveItemOne{}}{\numProbeCFourVOneRuleHostTrendCiNinetyFiveItemTwo{}}
(host-electronic) and \numProbeCFourVOneRuleSiteTrendRho{}
\ci{\numProbeCFourVOneRuleSiteTrendCiNinetyFiveItemOne{}}{\numProbeCFourVOneRuleSiteTrendCiNinetyFiveItemTwo{}}
(site-electronic). The trained probe exceeds the shuffled-label control at every budget the rule names, by
\numProbeCFourVOneRuleHostControlsSixHundredFiftyFourShuffledMean{}
\ci{\numProbeCFourVOneRuleHostControlsSixHundredFiftyFourShuffledCiNinetyFiveItemOne{}}{\numProbeCFourVOneRuleHostControlsSixHundredFiftyFourShuffledCiNinetyFiveItemTwo{}}
for the host class at \numCfgBudgetsHostsItemSix{} hosts. It does not exceed the untrained network at any of them.
At \numCfgBudgetsHostsItemFour{} hosts it lies below it (host-electronic
\numProbeCFourVOneRuleHostControlsTwoHundredInitZeroMean{}
\ci{\numProbeCFourVOneRuleHostControlsTwoHundredInitZeroCiNinetyFiveItemOne{}}{\numProbeCFourVOneRuleHostControlsTwoHundredInitZeroCiNinetyFiveItemTwo{}},
site-electronic \numProbeCFourVOneRuleSiteControlsTwoHundredInitZeroMean{}
\ci{\numProbeCFourVOneRuleSiteControlsTwoHundredInitZeroCiNinetyFiveItemOne{}}{\numProbeCFourVOneRuleSiteControlsTwoHundredInitZeroCiNinetyFiveItemTwo{}}),
and at \numCfgBudgetsHostsItemSix{} hosts the differences are
\numProbeCFourVOneRuleHostControlsSixHundredFiftyFourInitZeroMean{}
\ci{\numProbeCFourVOneRuleHostControlsSixHundredFiftyFourInitZeroCiNinetyFiveItemOne{}}{\numProbeCFourVOneRuleHostControlsSixHundredFiftyFourInitZeroCiNinetyFiveItemTwo{}}
and \numProbeCFourVOneRuleSiteControlsSixHundredFiftyFourInitZeroMean{}
\ci{\numProbeCFourVOneRuleSiteControlsSixHundredFiftyFourInitZeroCiNinetyFiveItemOne{}}{\numProbeCFourVOneRuleSiteControlsSixHundredFiftyFourInitZeroCiNinetyFiveItemTwo{}},
intervals that include zero. C4 therefore does not hold, and under the registered rule the probe is reported here
rather than in the article (Figure~\ref{fig:probe}). The collaborator's probe on the primary network (v2, post hoc,
descriptive) shows the same pattern, with trend correlations of \numProbeDescriptiveSVTwoHostRuleTrendRho{} and
\numProbeDescriptiveSVTwoSiteRuleTrendRho{}, and C4 would not hold on v2 either. What is not established is that
training on formation energies adds linearly recoverable electronic information beyond what an untrained encoder of
the same structure already exposes. One reading, not tested here, is that the host-electronic descriptors are
largely functions of composition and structure, which an untrained encoder with element embeddings and distance
expansions already represents.

A sanity check confirms that the probe reads the vector that the network's final step reads. Applying the stored
readout of the checkpoint to the stored vectors reproduces the network's own test predictions to within
\numProbeFitVOneSanityZeroSixHundredFiftyFourHeadReplayMaxAbsDiffEV{}~eV at \numCfgBudgetsHostsItemSix{} hosts, the
level of single-precision arithmetic. A linear map from the same vectors to the network's own prediction, however,
recovers only part of it out of fold (\Rsq{} = \numProbeFitVOneSanityZeroSixHundredFiftyFourProbeToOwnPredictionRTwoOutOfFold{}
at \numCfgBudgetsHostsItemSix{} hosts, against \numProbeFitVOneSanityZeroTwentyFiveProbeToOwnPredictionRTwoOutOfFold{}
at \numCfgBudgetsHostsItemOne{}), because the readout is a nonlinear multilayer network and a linear probe need not
reproduce it. The probe therefore measures linear recoverability, a lower bound on what the vector encodes; this was
recorded as found and not adjusted (\texttt{results/probe\_fit\_v1.json}, key \texttt{payload.sanity}).

\begin{figure}[htbp]
\centering
\sifig{fig6_probe}
\caption{Latent probe: linear recovery of the DFT electronic descriptors from the learned description of the vacancy
site and its host in the structure-only network (registered analysis on the originally registered network, v1). A
linear map, fitted on the training hosts of each resample and budget, asks how well the learned vectors that the
network's final step reads (vacancy site, pooled host, global state) predict each descriptor given to D, scored by
\Rsq{} on the test sites. Trained probe: dashed line, open markers and hatched band (dashed bars in (c)). Controls in
grey: the same probe with descriptor rows permuted among training hosts (shuffled labels; squares) and the probe on
an untrained network of the same architecture (the registered random-initialisation control; triangles); two further
untrained networks (light grey, dotted, no band) show the spread of that control and are descriptive.
(\textbf{a}) Host-electronic and (\textbf{b}) site-electronic class-mean test \Rsq{} against training hosts on a
logarithmic axis: means over the resamples with two-sided intervals of the resample bootstrap; the dotted line is
zero. (\textbf{c}) Test \Rsq{} per descriptor at the largest budget, host-electronic descriptors above the grey rule
and site-electronic below. The registered rule does not hold, so the probe is placed in the Supplementary Materials.
The collaborator's probe on the primary network (v2) is post hoc and is not drawn. Records:
\texttt{results/probe\_c4\_v1.json}, \texttt{results/probe\_descriptive.json}.}
\label{fig:probe}
\end{figure}
\FloatBarrier

\section{Structure renders and the geometry of the vacancy site}
\label{si:structures}

The networks read the pristine host supercell of the reference calculation with the oxygen atom to be removed
marked, never the relaxed defect cell. The two renders of this work show that input for two of the screened
oxides, wurtzite ZnO (\artfigoverview{}b of the article) and rutile SnO$_2$
(Figure~\ref{fig:sno2}), and Table~\ref{tab:shell} lists the first cation shell of the removed oxygen in each, read
from the pristine supercells of the release. In wurtzite ZnO the oxygen is tetrahedrally coordinated by four Zn
atoms: three Zn--O bonds of \numRenderZnOFirstShellSeventyFiveDistanceA{}~\AA{} and one of
\numRenderZnOFirstShellZeroDistanceA{}~\AA{}. In rutile SnO$_2$ the oxygen sits in a
planar three-fold site with one Sn--O bond of \numRenderSnOTwoFirstShellZeroDistanceA{}~\AA{} and two of
\numRenderSnOTwoFirstShellEighteenDistanceA{}~\AA{}. Removing the oxygen breaks four and three cation--oxygen bonds respectively, and the
two electrons of the neutral vacancy are left in this cation environment. The
geometry that the network reads is that of the pristine host: in the release's relaxed neutral defect cell of ZnO,
no atom lies more than \numStructureAuditZnoPristineVsRelaxedDisplacementAMaxOfMatchedAtoms{}~\AA{} from its pristine
position, and that relaxation is an output of the defect calculation that the models are meant to replace.

Both renders use the same camera conventions, lights and materials. The SnO$_2$ render was made with Blender
Cycles on the CPU with \numRenderSnOTwoSamples{} samples and denoising, at \numRenderSnOTwoDpi{} dpi over the text
width. Its cell is shifted by lattice translations only, so that the vacancy lies near the faces closest to the
viewer, and the planar three-fold site is seen nearly edge-on; Figure~\ref{fig:sno2} draws the three Sn neighbours
enlarged (parameters in \texttt{figures/si/structure\_sno2\_vacancy.render.json}).

\begin{table}[htbp]
\caption{First cation shell of the removed oxygen in the two rendered hosts (pristine PBEsol(+$U$) supercells of
the release). Atoms: atoms of the supercell; vacancy atom: index of the removed oxygen in the supercell; bond
cutoff: the cation--oxygen distance below which a bond is drawn in the renders. Sources:
\texttt{figures/main/fig1\_zno\_vacancy.render.json} and \texttt{figures/si/structure\_sno2\_vacancy.render.json}
(\texttt{first\_shell}).}
\label{tab:shell}
\centering\small
% Generated by paper/si/tools/make_si_tables.py from figures/main/fig1_zno_vacancy.render.json, figures/si/structure_sno2_vacancy.render.json. Do not edit.
\begin{tabular}{@{}lrrlrr@{}}
\toprule
Host & Atoms & Vacancy atom & Neighbour & M--O distance (\AA) & Bond cutoff (\AA) \\
\midrule
wurtzite ZnO & \numRenderZnONAtoms{} & \numRenderZnOVacancyAtomIndex{} & Zn & \numRenderZnOFirstShellSeventyFiveDistanceA{} & \numRenderZnOCationOBondCutoffAZn{} \\
 &  &  & Zn & \numRenderZnOFirstShellSeventyEightDistanceA{} &  \\
 &  &  & Zn & \numRenderZnOFirstShellOneHundredThirtyFiveDistanceA{} &  \\
 &  &  & Zn & \numRenderZnOFirstShellZeroDistanceA{} &  \\
\addlinespace
rutile SnO$_2$ & \numRenderSnOTwoNAtoms{} & \numRenderSnOTwoVacancyAtomIndex{} & Sn & \numRenderSnOTwoFirstShellZeroDistanceA{} & \numRenderSnOTwoCationOBondCutoffASn{} \\
 &  &  & Sn & \numRenderSnOTwoFirstShellNineteenDistanceA{} &  \\
 &  &  & Sn & \numRenderSnOTwoFirstShellEighteenDistanceA{} &  \\
\bottomrule
\end{tabular}

\end{table}

\begin{figure}[htbp]
\centering
\sifig{structure_sno2_vacancy}
\caption{Pristine supercell of rutile SnO$_2$ from the database (host \mpsnotwo{}, \numRenderSnOTwoNAtoms{} atoms), with the
oxygen to be removed (site O1) and its three nearest Sn neighbours (Table~\ref{tab:shell}) joined into a triangle.
The amber sphere marks the removed oxygen, Sn is grey-blue and O red, and Sn--O bonds shorter than
\numRenderSnOTwoCationOBondCutoffASn{}~\AA{} are drawn except across the cell boundary. Data:
\texttt{figures/si/structure\_sno2\_vacancy.render.json}.}
\label{fig:sno2}
\end{figure}
\FloatBarrier

\section{Reproducibility: code, records, deposit and provenance}
\label{si:repro}

\subsection{Code release and run records}

The code is released at \url{https://github.com/pa1aash/dft-gnn} under the MIT licence. Annotated tags mark the
states of the repository that the results depend on (Table~\ref{tab:tags}); the results reported in the article are
frozen at \tagresults{}, and the final code tag of the release is \texttt{[FINAL-TAG-TBD]}, to be set when the
deposit is minted.

\begin{table}[htbp]
\caption{Annotated tags of the code repository and what each marks (\texttt{git tag -n1}).}
\label{tab:tags}
\centering\small
% Generated by paper/si/tools/make_si_tables.py from the annotated tags of the repository (git tag -n1). Do not edit.
\begin{tabular}{@{}>{\raggedright}p{0.3\linewidth}>{\raggedright\arraybackslash}p{0.64\linewidth}@{}}
\toprule
Tag & What it marks \\
\midrule
\texttt{prereg-v1} & analysis plan frozen (registration of v1) \\
\texttt{tuned-v1} & tuned values of the originally registered network frozen \\
\texttt{g1-v1} & primary analysis of the learning curves (G1) \\
\texttt{nrp-code-v1} & pinned code of the cluster stages; the run ids of every cluster spec are computed with this name \\
\texttt{inputs-v1} & same commit as nrp-code-v1; release holding the processed inputs, checkpoints and MACE weights \\
\texttt{ingest-v1} & collaborator cluster analysis merged \\
\texttt{nrp-code-v2} & bug-fix tag of nrp-code-v1 for the relax, embed and geomeval reruns; run ids keep nrp-code-v1 \\
\texttt{nrp-code-v3} & result records also store the matgl version (record keeping only; no cluster run uses it) \\
\texttt{results-v1} & results freeze; headline numbers in results/freeze\_table.json \\
\texttt{prov-tune-v2, prov-prod-v2, prov-ckpt-diag} & code commits of the collaborator's v2 tuning, v2 production and checkpoint-diagnostic records \\
\bottomrule
\end{tabular}

\end{table}

Every training or analysis run writes one JSON record through a single function, \texttt{write\_result}, from a
clean working tree; the record holds the full configuration, the commit of the code, the package versions, the
seeds, the device (and, for the cluster runs, the GPU model) and the time, and the record of a training run names the hash of its per-site prediction table
(since \tagnrpthree{} it also stores the matgl version). A training run is named by its run id, a prefix of the
\shatwo{} digest of a canonical description of the run: model, hyperparameters, split, resample, budget, seed, code
reference and the hash of the graph-store manifest, plus any non-default option such as an epoch override, a
descriptor ablation or the component runs of the staged model (\texttt{src/dftgnn/train/\_\_init\_\_.py},
\texttt{run\_id}). The same description therefore always gives the same id, and two runs with the same id are the
same computation. For the cluster stages the code reference in the id is the tag \tagnrpone{} instead of the
generating commit, so the ids computed on the cluster and on the workstation agree; a worker runs such a job only if
the tag is an ancestor of its checkout and the code, configuration, split and job-specification directories are
unchanged since the tag. The bug-fix tag \tagnrptwo{} keeps the run ids of \tagnrpone{}. Named analysis records
(for example \texttt{results/g1\_primary.json}) are read by \texttt{paper/tools/make\_numbers.py}, which turns every
number of the article and of these Supplementary Materials into a macro with its source pointer; no number in
either document is typed by hand.

\subsection{Deposit and data availability}

The data and models are deposited in Zenodo under \ccby{} (DOI \texttt{[ZENODO-DOI-TBD]}, to be minted at
submission; \texttt{deposit/README.md}). The record holds the processed site table of the universe, the graph store
of the pristine supercells with its manifest, the release supercells of the hosts as released, the unit cells and
supercell tiling used for the machine-learned relaxations, all splits, the job specifications of the cluster stages,
the per-site prediction table and the record of every run, the checkpoints trained before the cluster stages, the
MACE-MP-0 weights used for the relaxations, and a \shatwo{} manifest of every file. Checkpoints, relaxed structures and
embeddings of the cluster runs are too large for the record; they are listed by \shatwo{} and are available from the
corresponding author on request. All formation energies, descriptors and structures derive from the oxygen-vacancy
database of Kumagai et al.\ \citep{Kumagai2021insights}, used at commit \releasecommit{} under \ccby{}; no new DFT
calculation was made.

\subsection{Provenance labels}

Each result group carries one of three labels, assigned by chronology rather than by outcome
(Table~\ref{tab:provenance}). The originally registered network (v1) and its rules were fixed in the analysis plan
at \tagprereg{} (\dateprereg{}) and its tuned values at \tagtuned{} (\datetunedvone{}). The primary network (v2) was
proposed after test-set diagnostics of v1 (design rule \datevtwoprotocol{}), selected on validation data only
(\datevtwoselect{}), retuned and frozen (\datevtwofreeze{}), and designated primary by author decision on
\dateauthordecision{}; every v2 result is post hoc and none is described as registered.

{\small
% Generated by paper/si/tools/make_si_tables.py from docs/provenance_labels.md. Do not edit.
\begin{longtable}{@{}r>{\raggedright}p{0.36\linewidth}>{\raggedright}p{0.2\linewidth}>{\raggedright\arraybackslash}p{0.3\linewidth}@{}}
\caption{Provenance label of every result group, transcribed from \texttt{docs/provenance\_labels.md} (groups 1 to 19: the collaborator analysis at its ingest; groups 20 to 35: the analyses added at the results freeze). PRE-REGISTERED: the originally registered backbone (v1) with the deciding rule committed before the results; POST HOC: the primary backbone (v2), or a rule committed after the results it decides; EXTENSION: compute or analysis beyond the registered plan. The labels record chronology, not quality. The records, rule commits and first-result times of every group are in the same file.}\label{tab:provenance}\\
\toprule
\# & Group & Network & Label \\
\midrule
\endfirsthead
\toprule
\# & Group & Network & Label \\
\midrule
\endhead
\bottomrule
\endlastfoot
1 & G1 primary (C1, C2), pre-registered arm & v1 & PRE-REGISTERED \\
2 & G1 on the v2 arm & v2 & POST HOC \\
3 & LOCO, S and D (\S{}11) & v1 (a654 values) & PRE-REGISTERED \\
4 & LOCO baselines B0 and RF-Kumagai on the folds & none & EXTENSION \\
5 & LOCO on the v2 arm & v2 & POST HOC \\
6 & staging (C3a): P1-v2 training and P-v2 evaluation & v2 (P1-v2 = \texttt{tuned-v1} P1 values + \texttt{init: kaiming}, not re-tuned) & POST HOC \\
7 & geometry evaluation (C3b) & v2 (S-v2, D-v2, P-v2 at B = 654 and 200) & POST HOC \\
8 & MLIP relaxations & none & PRE-REGISTERED \\
9 & latent probe (C4) & v2 & POST HOC \\
10 & screening (C5), pre-registered S & v1 (sweep S at B = 654) & PRE-REGISTERED \\
11 & screening (C5), S-v2 column & v2 & POST HOC \\
12 & descriptor ablations and permutation importances & v1 & EXTENSION \\
13 & D-late trained in the sweep protocol & v1 (D-late \texttt{tuned-v1} values) & EXTENSION \\
14 & budget $\times$ tuned-setting crossing & v1 & EXTENSION \\
15 & checkpoint gradient and flag diagnostics & v1 (retrained diagnostic checkpoints) & EXTENSION \\
16 & C0 tuning-overlap check & v1 (official C0 predictions) & EXTENSION \\
17 & v2 design screen & v2 candidates & POST HOC \\
18 & v2 retuning & v2 & POST HOC \\
19 & v2 production & v2 (\texttt{configs/\allowbreak{}tuned\_\allowbreak{}v2.yaml}) & POST HOC \\
20 & staging (C3a), registered analysis, step A2 & v1 & PRE-REGISTERED; the four-outcome TOST rule that decides its outcome is the collaborator rule of 2026-10-09, ratified 2026-10-11, POST HOC (as in \#6) \\
21 & staging (C3a-v2), re-analysis of \#6 with the G1 estimator, step A2 & v2 & POST HOC \\
22 & geometry evaluation (C3b) on our 818 relaxations at B = 654, with the non-converged sensitivity, step A3 & v1 & PRE-REGISTERED \\
23 & geometry at B = 200 & v1 & DESCRIPTIVE EXTRA, not pre-registered (label of \texttt{docs/\allowbreak{}a3\_\allowbreak{}notes.md}, kept) \\
24 & geometry: TOST of $\Delta$MAE\_P - $\Delta$MAE\_S and slope differences & v1 & POST HOC \\
25 & geometry: cross-check against the collaborator's 809 relaxations and his v2 evaluation & his v2 (\texttt{results/\allowbreak{}mlip\_\allowbreak{}eval\_\allowbreak{}v2.json}, \#7) & POST HOC (descriptive) \\
26 & latent probe (C4), registered analysis, step A4 & v1 (60 seed-0 S checkpoints of the sweep) & PRE-REGISTERED \\
27 & probe: within-host R$^{2}$ and per-descriptor values (PRB-WH) & v1 & DESCRIPTIVE (not part of the C4 rule; label of \texttt{docs/\allowbreak{}claims.yaml}, kept) \\
28 & probe on S-v2 (C4-v2), recomputed from \#9 & v2 & POST HOC (descriptive) \\
29 & precision extension (N4), step A1 & v1 & EXTENSION \\
30 & precision extension (N4), v2 & v2 & POST HOC (v2 backbone); EXTENSION in compute \\
31 & grouped K-fold CV (N3), step A1 & v1 & EXTENSION \\
32 & grouped K-fold CV (N3), v2 & v2 & POST HOC (v2 backbone); EXTENSION in compute \\
33 & high-moment sensitivity (SENS1), step A5 & v1 & EXTENSION \\
34 & CGCNN cross-check (SENS2), untuned, descriptive, step A5 & its own architecture; MEGNet v1 and v2 beside it & EXTENSION \\
35 & capped 600-epoch twins (SENS3), step A5 & v1 & POST HOC \\
\end{longtable}
}

\subsection{Deviations from the registered plan}

Every departure from the analysis plan is logged in \texttt{docs/deviations.md} before the run or analysis it
affects, with the plan section, the reason and the impact, and no row is edited after it is written; results
generated before a deviation are kept and reported. The rows fall into six kinds. Clarifications fix what the plan
left open without changing a registered decision: the early-stopping threshold, the seed rules, the tuning protocol,
the bootstrap and the rule for the smallest qualifying budget, the definitions of the geometry, probe and
leave-one-chemistry-out analyses, and operational details of the job queue. Deviations add an arm or an analysis
without removing one: the unselected injection route as a sensitivity, and the v2 network with its design screen,
selection rule and use in the staging, geometry and probe analyses. Extensions add compute or analysis beyond the
plan, each logged before its first result: the precision resamples, the grouped cross-validation, the within-host
breakdown, the rerun of the relaxations with an extended tiling rule, the repeatability study and the analyses of
the results freeze. One correction restricts the statement that the training cap never binds to the registered
check of the cap, because runs of the registered sweep did stop at the cap. One post hoc ratification records the
author sign-off, missing from the repository, of three collaborator rows: the equivalence rule and four-outcome
classification of the staging and geometry comparisons, the reporting rules for the smallest qualifying budget, and
the use of v2 for the staging, geometry and probe analyses; all three stay post hoc. Two author decisions designate v2 as the primary network and freeze the results at \tagresults{}; after the
freeze no experiment, run or analysis was added. The full log, with the commit time of every row, is part of the
code release.


\FloatBarrier
\bibliographystyle{mdpi}
\bibliography{../refs}

\end{document}
